\section{AIME II 2018}
        \begin{enumerate}
        	\item (\textit{AIME II 2018}) Points $A$, $B$, and $C$ lie in that order along a straight path where the distance from $A$ to $C$ is $1800$ meters. Ina runs twice as fast as Eve, and Paul runs twice as fast as Ina. The three runners start running at the same time with Ina starting at $A$ and running toward $C$, Paul starting at $B$ and running toward $C$, and Eve starting at $C$ and running toward $A$. When Paul meets Eve, he turns around and runs toward $A$. Paul and Ina both arrive at $B$ at the same time. Find the number of meters from $A$ to $B$.


 
 \textbf{Answer: }800



\textbf{Solution 1}

We know that in the same amount of time given, Ina will run twice the distance of Eve, and Paul would run quadruple the distance of Eve. Let's consider the time it takes for Paul to meet Eve: Paul would've run 4 times the distance of Eve, which we can denote as $d$. Thus, the distance between $B$ and $C$ is $4d+d=5d$. In that given time, Ina would've run twice the distance $d$, or $2d$ units. 

 Now, when Paul starts running back towards $A$, the same amount of time would pass since he will meet Ina at his starting point. Thus, we know that he travels another $4d$ units and Ina travels another $2d$ units. 

 Therefore, drawing out the diagram, we find that $2d+2d+4d+d=9d=1800 \implies d=200$, and distance between $A$ and $B$ is the distance Ina traveled, or $4d=4(200)=\boxed{800}$

 



\textbf{Solution 2}

Let $x$ be the distance from $A$ to $B$. Then the distance from $B$ to $C$ is $1800-x$. Since Eve is the slowest, we can call her speed $v$, so that Ina's speed is $2v$ and Paul's speed is $4v$. 

 For Paul and Eve to meet, they must cover a total distance of $1800-x$ which takes them a time of $\frac{1800-x}{4v+v}$. Paul must run the same distance back to $B$, so his total time is $\frac{2(1800-x)}{5v}$. 

 For Ina to reach $B$, she must run a distance of $x$ at a speed of $2v$, taking her a time of $\frac{x}{2v}$. 

 Since Paul and Ina reach $B$ at the same time, we know that $\frac{2(1800-x)}{5v} = \frac{x}{2v}$ (notice that $v$ cancels out on both sides). Solving the equation gives $x = \boxed{800}$.

 


 


	\item (\textit{AIME II 2018}) Let $a_{0} = 2$, $a_{1} = 5$, and $a_{2} = 8$, and for $n > 2$ define $a_{n}$ recursively to be the remainder when $4(a_{n-1} + a_{n-2} + a_{n-3})$ is divided by $11$. Find $a_{2018} \cdot a_{2020} \cdot a_{2022}$.


 
 \textbf{Answer: }112



\textbf{Solution 1}

When given a sequence problem, one good thing to do is to check if the sequence repeats itself or if there is a pattern.

 After computing more values of the sequence, it can be observed that the sequence repeats itself every 10 terms starting at $a_{0}$.

 $a_{0} = 2$,$a_{1} = 5$,$a_{2} = 8$,$a_{3} = 5$,$a_{4} = 6$,$a_{5} = 10$,$a_{6} = 7$,$a_{7} = 4$,$a_{8} = 7$,$a_{9} = 6$,$a_{10} = 2$,$a_{11} = 5$,$a_{12} = 8$,$a_{13} = 5$

 We can simplify the expression we need to solve to $a_{8}\cdot a_{10} \cdot a_{2}$.

 Our answer is $7 \cdot 2 \cdot 8$ $= \boxed{112}$.

 



\textbf{Solution 2 (Overkill, but advantageous if numbers are bigger)}

Notice that the characteristic polynomial of this is $x^3-4x^2-4x-4\equiv 0\pmod{11}$

 Then since $x\equiv1$ is a root, using Vieta's formula, the other two roots $r,s$ satisfy $r+s\equiv3$ and $rs\equiv4$. 

 Let $r=7+d$ and $s=7-d$.

 We have $49-d^2\equiv4$ so $d\equiv1$. We found that the three roots of the characteristic polynomial are $1,6,8$.

 Now we want to express $a_n$ in an explicit form as $a(1^n)+b(6^n)+c(8^n)\pmod{11}$.

 Plugging in $n=0,1,2$ we get

 $(*)$$a+b+c\equiv2,$

 $(**)$$a+6b+8c\equiv5,$

 $(***)$$a+3b+9c\equiv8$

 $\frac{(***)-(*)}{2}$$\implies b+4c\equiv3$ and $(***)-(**)$$\implies -3b+c\equiv3$

 so $a\equiv6,$ $b\equiv1,$ and $c\equiv6$

 Hence, $a_n\equiv 6+(6^n)+6(8^n)\equiv(2)^{-n\pmod{10}}+(2)^{3n-1\pmod{10}}-5\pmod{11}$

 Therefore

 $a_{2018}\equiv4+8-5=7$

 $a_{2020}\equiv1+6-5=2$

 $a_{2022}\equiv3+10-5=8$

 And the answer is $7\times2\times8=\boxed{112}$

 


	\item (\textit{AIME II 2018}) Find the sum of all positive integers $b < 1000$ such that the base-$b$ integer $36_{b}$ is a perfect square and the base-$b$ integer $27_{b}$ is a perfect cube.


 
 \textbf{Answer: }371



\textbf{Solution 1}

The first step is to convert $36_{b}$ and $27_{b}$ into base-10 numbers. Then, we can write 
 \[36_{b} = 3b + 6\] and 
 \[27_{b} = 2b + 7\]. It should also be noted that $8 \leq b < 1000$.

 Because there are less perfect cubes than perfect squares for the restriction we are given on $b$, it is best to list out all the perfect cubes. Since the maximum $b$ can be is 1000 and $2$ • $1000 + 7 = 2007$, we can list all the perfect cubes less than 2007.

 Now, $2b + 7$ must be one of 
 \[3^3, 4^3, ... , 12^3\]. However, $2b + 7$ will always be odd, so we can eliminate the cubes of the even numbers and change our list of potential cubes to 
 \[3^3, 5^3, 7^3, 9^3\text{, and }11^3\].

 Because $3b + 6$ is a perfect square and is clearly divisible by 3, it must be divisible by 9, so $b + 2$ is divisible by 3. Thus the cube, which is 
 \[2b + 7 = 2(b + 2) + 3\], must also be divisible by 3. Therefore, the only cubes that $2b + 7$ could potentially be now are $3^3$ and $9^3$.

 We need to test both of these cubes to make sure $3b + 6$ is a perfect square.

 $\textbf{Case 1:}$ If we set 
 \[3^3 = 2b + 7\]so 
 \[b = 10\]. If we plug this value of b into $3b + 6$, the expression equals $36$, which is indeed a perfect square.

 $\textbf{Case 2:}$ If we set 
 \[9^3 = 2b + 7\]so 
 \[b = 361\]. If we plug this value of b into $3b + 6$, the expression equals $1089$, which is $33^2$.

 We have proven that both $b = 10$ and $b = 361$ are the only solutions, so 
 \[10 + 361 = \boxed{371}\]

 



\textbf{Solution 2}

The conditions are:
 \[3b+6 = n^2\]
 \[2b+7 = m^3\]We can see $n$ is multiple is 3, so let $n=3k$, then $b= 3k^2-2$. Substitute $b$ into second condition and we get $m^3=3(2k^2+1)$. Now we know $m$ is both a multiple of 3 and odd. Also, $m$ must be smaller than 13 for $b$ to be smaller than 1000. So the only two possible values for $m$ are 3 and 9. Test and they both work. The final answer is $10 + 361 =$ $\boxed{371}$.-Mathdummy

 



\textbf{Solution 3}

As shown above, let 
 \[3b+6 = n^2\]
 \[2b+7 = m^3\] such that 
 \[6b+12=2n^2\] 
 \[6b+21=3m^3\]

 Subtracting the equations we have 
 \[3m^3-2n^2=9 \implies 3m^3-2n^2-9=0.\]

 We know that $m$ and $n$ both have to be integers, because then the base wouldn't be an integer. Furthermore, any integer solution $m$ must divide $9$ by the Rational Root Theorem.

 We can instantly know $m \neq -9,-3,-1,1$ since those will have negative solutions.

 When $m=3$ we have $n=6$, so then $b=10$

 When $m=9$ we have $n=33$, so then $b=361$

 Therefore, the sum of all possible values of $b$ is 
 \[10+361=\boxed{371}.\]

 


	\item (\textit{AIME II 2018}) In equiangular octagon $CAROLINE$, $CA = RO = LI = NE =$ $\sqrt{2}$ and $AR = OL = IN = EC = 1$. The self-intersecting octagon $CORNELIA$ encloses six non-overlapping triangular regions. Let $K$ be the area enclosed by $CORNELIA$, that is, the total area of the six triangular regions. Then $K = \frac{a}{b}$, where $a$ and $b$ are relatively prime positive integers. Find $a + b$.


 
 \textbf{Answer: }023



\textbf{Solution}

We can draw $CORNELIA$ and introduce some points.

 
                    \begin{figure}[H]
                        \centering
                        \includegraphics[width=0.5\textwidth]{Solution/AIME_II_2018_4_Solution_Figure_1.png}
                    \end{figure}
                    

 The diagram is essentially a 3x3 grid where each of the 9 squares making up the grid have a side length of 1.

 In order to find the area of $CORNELIA$, we need to find 4 times the area of $\bigtriangleup$$ACY$ and 2 times the area of $\bigtriangleup$$YZW$.

 Using similar triangles $\bigtriangleup$$ARW$ and $\bigtriangleup$$YZW$(We look at their heights), $YZ$ $=$ $\frac{1}{3}$. Therefore, the area of $\bigtriangleup$$YZW$ is $\frac{1}{3}\cdot\frac{1}{2}\cdot\frac{1}{2}$ $=$ $\frac{1}{12}$

 Since $YZ$ $=$ $\frac{1}{3}$ and $XY = ZQ$, $XY$ $=$ $\frac{1}{3}$ and $CY$ $=$ $\frac{4}{3}$.

 Therefore, the area of $\bigtriangleup$$ACY$ is $\frac{4}{3}\cdot$ $1$ $\cdot$ $\frac{1}{2}$ $=$ $\frac{2}{3}$

 Our final answer is $\frac{1}{12}$ $\cdot$ $2$ $+$ $\frac{2}{3}$ $\cdot$ $4$ $=$ $\frac{17}{6}$

 $17 + 6 =$ $\boxed{023}$

 



\textbf{Solution 2}

$CAROLINE$ is essentially a plus sign with side length 1 with a few diagonals, which motivates us to coordinate bash. We let $N = (1, 0)$ and $E = (0, 1)$. To find $CORNELIA$'s self intersections, we take 

 
 \[CO = y = 2, AI = y = -3x + 6, RN = y = 3x - 3\]

 And plug them in to get $C_1 = \left(\frac{4}{3}, 2 \right)$ where $C_1$ is the intersection of $CO$ and $AI$, and $C_2 = \left(\frac{5}{3}, 2 \right)$ is the intersection of $RN$ and $CO$.

 We also track the intersection of $AI$ and $RN$ to get $\left(\frac{3}{2}, \frac{3}{2} \right)$.

 By vertical symmetry, the other 2 points of intersection should have the same x-coordinates. We can then proceed with Solution 1 to calculate the area of the triangle (compare the $y$-coordinates of $A,R,I,N$ and $CO$ and $EL$).

 


	\item (\textit{AIME II 2018}) Suppose that $x$, $y$, and $z$ are complex numbers such that $xy = -80 - 320i$, $yz = 60$, and $zx = -96 + 24i$, where $i$ $=$ $\sqrt{-1}$. Then there are real numbers $a$ and $b$ such that $x + y + z = a + bi$. Find $a^2 + b^2$.


 
 \textbf{Answer: }074



\textbf{Solution 1 (Euler's formula and Substitution)}

The First (pun intended) thing to notice is that $xy$ and $zx$ have a similar structure (i.e. the real and imaginary parts differ by a factor of 4), but are not exactly conjugates. So let's take out the magnitudes of both, and see if we can multiply a root of unity to get the other. It turns out that root of unity is $e^{\frac{3\pi i}{2}}$. Anyway this results in getting that $\left(\frac{-3i}{10}\right)y=z$. 

 Then substitute this into $yz$ to get, after some calculation, that $y=10e^{\frac{5\pi i}{4}}\sqrt{2}$ and $z=-3e^{\frac{7\pi i}{4}}\sqrt{2}$. Then plug $z$ into $zx$, you could do the same thing with $xy$ but $zx$ looks like it's easier due to it being smaller. Anyway you get $x=20+12i$. Then add all three up, which it turns out easier than it seems because for $z$ and $y$ the $\sqrt{2}$ disappears after you expand the root of unity (e raised to a specific power). 

 Long story short, you get $x=20+12i, y=-3+3i, z=-10-10i \implies x+y+z=7+5i \implies a^2+b^2=\boxed{074}$.

 



\textbf{Solution 2}

First we evaluate the magnitudes. $|xy|=80\sqrt{17}$, $|yz|=60$, and $|zx|=24\sqrt{17}$. Therefore, $|x^2y^2z^2|=17\cdot80\cdot60\cdot24$, or $|xyz|=240\sqrt{34}$. Divide to find that $|z|=3\sqrt{2}$, $|x|=4\sqrt{34}$, and $|y|=10\sqrt{2}$.
                    \begin{figure}[H]
                        \centering
                        \includegraphics[width=0.5\textwidth]{Solution/AIME_II_2018_5_Solution_Figure_1.png}
                    \end{figure}
                    This allows us to see that the argument of $y$ is $\frac{\pi}{4}$, and the argument of $z$ is $-\frac{\pi}{4}$. We need to convert the polar form to a standard form. Simple trig identities show $y=10+10i$ and $z=3-3i$. More division is needed to find what $x$ is. 
 \[x=-20-12i\] 
 \[x+y+z=-7-5i\] 
 \[(-7)^2+(-5)^2=\boxed{74}\]
 \[QED\blacksquare\]Written by \href{https://artofproblemsolving.com/wiki/index.php?title=User:A1b2}{a1b2}

 



\textbf{Solution 3 (Pretty easy, no hard stuff, just watch ur arithmetic) :)}

Solve this system the way you would if the RHS of all equations were real. Multiply the first and 3rd equations out and then factor out $60$ to find $x^2$, then use standard techniques that are used to evaluate square roots of irrationals. let 
 \[x = c+di\], then you get 
 \[c^2 - d^2 = 256\] and 
 \[2cd = 480\] Solve to get $x$ as $20+12i$ and $-20-12i$. Both will give us the same answer, so use the positive one. Divide $-80-320i$ by  $x$, and you get $10+10i$ as $y$. This means that $z$ is a multiple of $1-i$ to get a real product, so you find $z$ is $3-3i$. Now, add the real and imaginary parts separately to get $-7-5i$, and calculate $a^2 + b^2$ to get $\boxed{074}$. ~minor latex improvements done by jske25 and jdong2006

 



\textbf{Solution 4}

Dividing the first equation by the second equation given, we find that $\frac{xy}{yz}=\frac{x}{z}=\frac{-80-320i}{60}=-\frac{4}{3}-\frac{16}{3}i \implies x=z\left(-\frac{4}{3}-\frac{16}{3}i\right)$. Substituting this into the third equation, we get $z^2=\frac{-96+24i}{-\frac{4}{3}-\frac{16}{3}i}=3\cdot \frac{-24+6i}{-1-4i}=3\cdot \frac{(-24+6i)(-1+4i)}{1+16}=3\cdot \frac{-102i}{17}=-18i$. Taking the square root of this is equivalent to halving the argument and taking the square root of the magnitude. Furthermore, the second equation given tells us that the argument of $y$ is the negative of that of $z$, and their magnitudes multiply to $60$. Thus, we have $z=\sqrt{-18i}=3-3i$ and $3\sqrt{2}\cdot |y|=60 \implies |y|=10\sqrt{2} \implies y=10+10i$. To find $x$, we can use the previous substitution we made to find that $x=z\left(-\frac{4}{3}-\frac{16}{3}i\right)=-\frac{4}{3}\cdot (3-3i)(1+4i)=-4(1-i)(1+4i)=-4(5+3i)=-20-12i$.Therefore, $x+y+z=(-20+10+3)+(-12+10-3)i=-7-5i \implies a^2+b^2=(-7)^2+(-5)^2=49+25=\boxed{074}$

 Solution by ktong

 



\textbf{Solution 5}

We are given that $xy=-80-320i$. Thus $y=\frac{-80-320i}{x}$. We are also given that $xz= -96+24i$. Thus $z=\frac{-96+24i}{x}$. We are also given that $yz$ = $60$. Substitute $y=\frac{-80-320i}{x}$ and $z=\frac{-96+24i}{x}$ into $yz$ = $60$. We have $\frac{(-80-320i)(-96+24i)}{x^2}=60$. Multiplying out $(-80-320i)(-96+24i)$ we get $(1920)(8+15i)$. Thus  $\frac{1920(8+15i)}{x^2} =60$. Simplifying this fraction we get $\frac{32(8+15i)}{x^2}=1$. Cross-multiplying the fractions we get $x^2=32(8+15i)$ or $x^2= 256+480i$. Now we can rewrite this as $x^2-256=480i$. Let $x= (a+bi)$.Thus $x^2=(a+bi)^2$ or $a^2+2abi-b^2$. We can see that $a^2+2abi-b^2-256=480i$ and thus $2abi=480i$ or $ab=240$.We also can see that $a^2-b^2-256=0$ because there is no real term in $480i$. Thus $a^2-b^2=256$ or $(a+b)(a-b)=256$. Using the two equations $ab=240$ and $(a+b)(a-b)=256$ we solve by doing system of equations that $a=-20$ and $b=-12$. And $x=a+bi$ so $x=-20-12i$. Because $y=\frac{-80-320i}{x}$, then $y=\frac{-80-320i}{-20-12i}$. Simplifying this fraction we get $y=\frac{-80(1+4i)}{-4(5+3i)}$ or $y=\frac{20(1+4i)}{(5+3i)}$. Multiplying by the conjugate of the denominator ($5-3i$) in the numerator and the denominator and  we get $y=\frac{20(17-17i)}{34}$. Simplifying this fraction we get $y=10-10i$. Given that $yz$ = $60$ we can substitute $(10-10i)(z)=60$ We can solve for z and get $z=3+3i$. Now we know what $x$, $y$, and $z$ are, so all we have to do is plug and chug. $x+y+z= (-20-12i)+(10+10i)+(3-3i)$ or $x+y+z= -7-5i$ Now $a^2 +b^2=(-7)^2+(-5)^2$ or $a^2 +b^2 = 74$. Thus $74$ is our final answer.(David Camacho)

 



\textbf{Solution 6}

We observe that by multiplying $xy,$ $yz,$ and $zx,$ we get $(xyz)^2=(-80-320i)(60)(-96+24i).$ Next, we divide $(xyz)^2$ by $(yz)^2$ to get $x^2.$ We have $x^2=\frac{(-80-320i)(60)(-96+24i)}{3600}=256+480i.$ We can write $x$ in the form of $a+bi,$ so we get $(a+bi)^2=256+480i.$ Then, $a^2-b^2+2abi=256+480i,$ $a^2-b^2=256,$ and $2ab=480.$ Solving this system of equations is relatively simple. We have two cases, $a=20, b=12,$ and $a=-20, b=-12.$ Case 1: $a=20, b=12,$ so $x=20+12i.$ We solve for $y$ and $z$ by plugging in $x$ to the two equations. We see$y=\frac{-80-320i}{20+12i}=-10-10i$ and $z=\frac{-96+24i}{20+12i}=-3+3i.$ $x+y+z=7+5i,$ so $a=7$ and $b=5.$ Solving, we end up with $7^2+5^2=\boxed{074}$ as our answer. Case 2: $a=-20, b=-12,$ so $x=-20-12i.$ Again, we solve for $y$ and $z.$ We find $y=\frac{-80-320i}{-20-12i}=10+10i,$ $z=\frac{-96+24i}{-20-12i}=3-3i,$ so $x+y+z=-7-5i.$ We again have $(-7)^2+(-5)^2=\boxed{074}.$Solution by Airplane50

 



\textbf{Solution 7 (Based on advanced mathematical knowledge)}

According to the Euler's Theory, we can rewrite $x$, $y$ and $z$ as 
 \[x=r_{1}e^{i{\theta}_1}\] 
 \[y=r_{2}e^{i{\theta}_2}\] 
 \[x=r_{3}e^{i{\theta}_3}\] As a result, 
 \[|xy|=r_{1}r_{2}=\sqrt{80^2+320^2}=80\sqrt{17}\] 
 \[|yz|=r_{2}r_{3}=60\] 
 \[|xz|=r_{1}r_{3}=\sqrt{96^2+24^2}=24\sqrt{17}\] Also, it is clear that 
 \[yz=r_{2}e^{i{\theta}_2}r_{3}e^{i{\theta}_3}=|yz|e^{i({\theta}_2+{\theta}_3)}=|yz|=60\] So ${\theta}_2+{\theta}_3=0$, or 
 \[{\theta}_2=-{\theta}_3\] Also, we have 
 \[xy=-80\sqrt{17}e^{i\arctan{4}}\] 
 \[yz=60\] 
 \[xz=-24\sqrt{17}e^{i\arctan{-\frac{1}{4}}}\] So now we have $r_{1}r_{2}=80\sqrt{17}$, $r_{2}r_{3}=60$, $r_{1}r_{3}=24\sqrt{17}$, ${\theta}_1+{\theta}_2=\arctan{4}$ and ${\theta}_1-{\theta}_2=\arctan {-\frac{1}{4}}$. Solve these above, we get 
 \[r_{1}=4\sqrt{34}\] 
 \[r_{2}=10\sqrt{2}\] 
 \[r_{3}=3\sqrt{2}\] 
 \[{\theta}_2=\frac{\arctan{4}-\arctan{-\frac{1}{4}}}{2}=\frac{\frac{\pi}{2}}{2}=\frac{\pi}{4}\] So we can get 
 \[y=r_{2}e^{i{\theta}_2}=10\sqrt{2}e^{i\frac{\pi}{4}}=10+10i\] 
 \[z=r_{3}e^{i{\theta}_3}=r_{3}e^{-i{\theta}_2}=3\sqrt{2}e^{-i\frac{\pi}{4}}=3-3i\] Use $xy=-80-320i$ we can find that 
 \[x=-20-12i\] So 
 \[x+y+z=-20-12i+10+10i+3-3i=-7-5i\] So we have $a=-7$ and $b=-5$.

 As a result, we finally get 
 \[a^2+b^2=(-7)^2+(-5)^2=\boxed{074}\]

 



\textbf{Solution 8}

We can turn the expression $x+y+z$ into $\sqrt{x^2+y^2+z^2+2xy+2yz+2xz}$, and this would allow us to plug in the values after some computations.

 Based off of the given products, we have

 $xy^2z=60(-80-320i)$

 $xyz^2=60(-96+24i)$

 $x^2yz=(-96+24i)(-80-320i)$.

 Dividing by the given products, we have

 $y^2=\frac{60(-80-320i)}{-96+24i}$

 $z^2=\frac{60(-96+24i)}{-80-320i}$

 $x^2=\frac{(-96+24i)(-80-320i)}{60}$.

 Simplifying, we get that this expression becomes $\sqrt{24+70i}$. This equals $\pm{(7+5i)}$, so the answer is $7^2+5^2=\boxed{74}$.

 $\textbf{-RootThreeOverTwo}$

 



\textbf{Solution 9 (A Little Rigorous, but Straightforward and Easy)}

Multiplying $xy \cdot yz \cdot zx = (xyz)^2$ we obtain $60 \cdot 960(16+30i)$ (too lazy to do $60 \cdot 960$, you don't need to). Taking the square root, we get $240\sqrt{16+30i}$. Letting $(a+bi)^2=16+30i,$ we have $a^2+2abi-b^2=16+30i.$ Thus, $(a+b)(a-b)=16,$ and $2ab=30.$ Guessing and checking, we get $a+bi=5+3i$. Therefore, $xyz=240(5+3i).$ Dividing this by each of the equations provided in the original problem, we get $x=20+12i,y=-10-10i,$ and $z=-3+3i$. $20+12i-10-10i-3+3i=7+5i$. Finally, $7^2+5^2=\boxed{074}.$

 


	\item (\textit{AIME II 2018}) A real number $a$ is chosen randomly and uniformly from the interval $[-20, 18]$. The probability that the roots of the polynomial


 
 \[x^4 + 2ax^3 + (2a - 2)x^2 + (-4a + 3)x - 2\]


 are all real can be written in the form $\dfrac{m}{n}$, where $m$ and $n$ are relatively prime positive integers. Find $m + n$.


 
 \textbf{Answer: }037



\textbf{Solution}

The polynomial we are given is rather complicated, so let's use \href{https://artofproblemsolving.com/wiki/index.php?title=Rational\_Root\_Theorem}{Rational Root Theorem} to see if it has any retinal roots. By Rational Root Theorem, $x = 1, -1, 2, -2$ are all possible rational roots. Upon plugging these roots into the polynomial, $x = -2$ and $x = 1$ make the polynomial equal 0 and thus, they are roots that we can factor out.

 The polynomial becomes:

 $(x - 1)(x + 2)(x^2 + (2a - 1)x + 1)$

 Since we know $1$ and $-2$ are real numbers, we only need to focus on the quadratic.

 Set the discriminant of the quadratic greater than or equal to 0, to ensure the remaining roots are real.

 $(2a - 1)^2 - 4 \geq 0$.

 This simplifies to:

 $a \geq \dfrac{3}{2}$

 or

 $a \leq -\dfrac{1}{2}$

 This means that the interval $\left(-\dfrac{1}{2}, \dfrac{3}{2}\right)$ is the "bad" interval. The length of the interval where $a$ can be chosen from is 38 units long, while the bad interval is 2 units long. Therefore, the "good" interval is 36 units long.

 $\dfrac{36}{38} = \dfrac{18}{19}$

 $18 + 19 = \boxed{037}$

 



\textbf{Solution 2}

This solution is essentially the same as solution 1, but it offers a different approach to factoring the polynomial. 

 We want to isolate $a$, so we rewrite $x^4 + 2ax^3 + (2a - 2)x^2 + (-4a + 3)x - 2$ as $x^4-2x^2+3x-2+2ax^3+2ax^2-4ax$.

 We see that $1$ and $-2$ are roots of $x^4-2x^2+3x-2$, so we can factor, yielding $(x-1)(x+2)(x^2-x+1)$.

 We can factor out $2ax$ from the other part of the polynomial, yielding $2ax(x^2+x-2)=2ax(x-1)(x+2)$.

 Thus, 

 \begin{align}x^4 + 2ax^3 + (2a - 2)x^2 + (-4a + 3)x - 2 &= x^4-2x^2+3x-2+2ax^3+2ax^2-4ax \\ &= (x-1)(x+2)(x^2-x+1) + 2ax(x-1)(x+2)\\&= (x-1)(x+2)(x^2+(2a-1)x+1). \\\end{align}

 The rest of the steps are the same as in Solution 1.

 


	\item (\textit{AIME II 2018}) Triangle $ABC$ has side lengths $AB = 9$, $BC =$ $5\sqrt{3}$, and $AC = 12$. Points $A = P_{0}, P_{1}, P_{2}, ... , P_{2450} = B$ are on segment $\overline{AB}$ with $P_{k}$ between $P_{k-1}$ and $P_{k+1}$ for $k = 1, 2, ..., 2449$, and points $A = Q_{0}, Q_{1}, Q_{2}, ... , Q_{2450} = C$ are on segment $\overline{AC}$ with $Q_{k}$ between $Q_{k-1}$ and $Q_{k+1}$ for $k = 1, 2, ..., 2449$. Furthermore, each segment $\overline{P_{k}Q_{k}}$, $k = 1, 2, ..., 2449$, is parallel to $\overline{BC}$. The segments cut the triangle into $2450$ regions, consisting of $2449$ trapezoids and $1$ triangle. Each of the $2450$ regions has the same area. Find the number of segments $\overline{P_{k}Q_{k}}$, $k = 1, 2, ..., 2450$, that have rational length.


 
 \textbf{Answer: }020



\textbf{Solution 1}

For each $k$ between $2$ and $2450$, the area of the trapezoid with $\overline{P_kQ_k}$ as its bottom base is the difference between the areas of two triangles, both similar to $\triangle{ABC}$. Let $d_k$ be the length of segment $\overline{P_kQ_k}$. The area of the trapezoid with bases $\overline{P_{k-1}Q_{k-1}}$ and $P_kQ_k$ is $\left(\frac{d_k}{5\sqrt{3}}\right)^2 - \left(\frac{d_{k-1}}{5\sqrt{3}}\right)^2 = \frac{d_k^2-d_{k-1}^2}{75}$ times the area of $\triangle{ABC}$. (This logic also applies to the topmost triangle if we notice that $d_0 = 0$.) However, we also know that the area of each shape is $\frac{1}{2450}$ times the area of $\triangle{ABC}$. We then have $\frac{d_k^2-d_{k-1}^2}{75} = \frac{1}{2450}$. Simplifying, $d_k^2-d_{k-1}^2 = \frac{3}{98}$. However, we know that $d_0^2 = 0$, so $d_1^2 = \frac{3}{98}$, and in general, $d_k^2 = \frac{3k}{98}$ and $d_k = \frac{\sqrt{\frac{3k}{2}}}{7}$. The smallest $k$ that gives a rational $d_k$ is $6$, so $d_k$ is rational if and only if $k = 6n^2$ for some integer $n$.The largest $n$ such that $6n^2$ is less than $2450$ is $20$, so $k$ has $\boxed{020}$ possible values.

 Solution by zeroman

 



\textbf{Solution 2}

We have that there are $2449$ trapezoids and $1$ triangle of equal area, with that one triangle being $AP_1Q_1$. Notice, if we "stack" the trapezoids on top of $\bigtriangleup AP_1Q_1$ the way they already are, we'd create a similar triangle, all of which are similar to $\bigtriangleup ABC$, and since the trapezoids and $\bigtriangleup AP_1Q_1$ have equal area, each of these similar triangles $AP_kQ_k$ have area $\frac{k}{2450}\left[ ABC\right]$, and so $\frac{\left[ AP_kQ_k\right]}{\left[ABC\right]}=\frac{k}{2450}$. We want the ratio of the side lengths $P_kQ_k:BC$. Since area is a 2-dimensional unit of measurement, and side lengths are 1-dimensional, the ratio is simply the square root of the areas, or 
 \[\frac{P_kQ_k}{BC}=\sqrt{\frac{k}{2450}}\] 
 \[\implies P_kQ_k=BC\cdot \sqrt{\frac{k}{2450}}=5\sqrt{3}\cdot\sqrt{\frac{k}{2450}}=\frac{1}{7}\cdot \sqrt{\frac{3k}{2}}=\frac{3}{7}\sqrt{\frac{k}{6}}\] 
 \[\implies k=6n^2<2450\] 
 \[\implies 0<n\leq 20\] so there are $\boxed{020}$ solutions. 

 



\textbf{Solution 3}

Let $T_1$ stand for $AP_1Q_1$, and $T_k = AP_kQ_k$. All triangles $T$ are similar by AA. Let the area of $T_1$ be $x$. The next trapezoid will also have an area of $x$, as given. Therefore, $T_k$ has an area of $kx$. The ratio of the areas is equal to the square of the scale factor for any plane figure and its image. Therefore, $P_k Q_k=P_1 Q_1\cdot \sqrt{k}$, and the same if $Q$ is substituted for $P$ throughout. We want the side $P_k Q_k$ to be rational. Setting up proportions: 
 \[5\sqrt{3} : \sqrt{2450}=35\sqrt{2}\] 
 \[\sqrt{6} : 14\] which shows that $P_1 Q_1=\frac{\sqrt{6}}{14}$. In order for $\sqrt{k} P_1 Q_1$ to be rational, $\sqrt{k}$ must be some rational multiple of $\sqrt{6}$. This is achieved at $\sqrt{k}=\sqrt{6}, 2\sqrt{6}, \ldots, 20\sqrt{6}$. We end there as $21\sqrt{6}=\sqrt{2646}$. There are 20 numbers from 1 to 20, so there are $\boxed{020}$ solutions.

 Solution by a1b2

 



\textbf{Solution 4: I didn't even qualify and I solved this}

We can also use area ratios to get that $P_nQ_n$ is $\sqrt(n)(P_1Q_1)$. So $P_{2450}Q_{2450} = 35 \sqrt(3) = 35\sqrt(2)P_1Q_1$ so $P_1Q_1$(call this $L_1$) is $\sqrt(6)/14$. So We want $L_1\sqrt(n)$ to be rational, or that $\sqrt(n) = j\sqrt(6)$. So $j$ can range from 1 to 20, so we get $\boxed{020}$.

 


	\item (\textit{AIME II 2018}) A frog is positioned at the origin of the coordinate plane. From the point $(x, y)$, the frog can jump to any of the points $(x + 1, y)$, $(x + 2, y)$, $(x, y + 1)$, or $(x, y + 2)$. Find the number of distinct sequences of jumps in which the frog begins at $(0, 0)$ and ends at $(4, 4)$.


 
 \textbf{Answer: }556



\textbf{Solution 1}

We solve this problem by working backwards. Notice, the only points the frog can be on to jump to $(4,4)$ in one move are $(2,4),(3,4),(4,2),$ and $(4,3)$. This applies to any other point, thus we can work our way from $(0,0)$ to $(4,4)$, recording down the number of ways to get to each point recursively. 

 $(0,0): 1$

 $(1,0)=(0,1)=1$

 $(2,0)=(0, 2)=2$

 $(3,0)=(0, 3)=3$

 $(4,0)=(0, 4)=5$

 $(1,1)=2$, $(1,2)=(2,1)=5$, $(1,3)=(3,1)=10$, $(1,4)=(4,1)= 20$

 $(2,2)=14, (2,3)=(3,2)=32, (2,4)=(4,2)=71$

 

 


	\item (\textit{AIME II 2018}) Octagon $ABCDEFGH$ with side lengths $AB = CD = EF = GH = 10$ and $BC = DE = FG = HA = 11$ is formed by removing 6-8-10 triangles from the corners of a $23$ $\times$ $27$ rectangle with side $\overline{AH}$ on a short side of the rectangle, as shown. Let $J$ be the midpoint of $\overline{AH}$, and partition the octagon into 7 triangles by drawing segments $\overline{JB}$, $\overline{JC}$, $\overline{JD}$, $\overline{JE}$, $\overline{JF}$, and $\overline{JG}$. Find the area of the convex polygon whose vertices are the centroids of these 7 triangles.


 
                    \begin{figure}[H]
                        \centering
                        \includegraphics[width=0.5\textwidth]{Problem/AIME_II_2018_9_Problem_Figure_1.png}
                    \end{figure}
                    


 
 \textbf{Answer: }184



\textbf{Solution 1 (Massive Shoelace)}

We represent octagon $ABCDEFGH$ in the coordinate plane with the upper left corner of the rectangle being the origin. Then it follows that $A=(0,-6), B=(8, 0), C=(19,0), D=(27, -6), E=(27, -17), F=(19, -23), G=(8, -23), H=(0, -17), J=(0, -\frac{23}{2})$. Recall that the centroid is $\frac{1}{3}$ way up each median in the triangle. Thus, we can find the centroids easily by finding the midpoint of the side opposite of point $J$. Furthermore, we can take advantage of the reflective symmetry across the line parallel to $BC$ going through $J$ by dealing with less coordinates and ommiting the $\frac{1}{2}$ in the shoelace formula. 

 By doing some basic algebra, we find that the coordinates of the centroids of $\bigtriangleup JAB, \bigtriangleup JBC, \bigtriangleup JCD, \bigtriangleup JDE$ are $\left(\frac{8}{3}, -\frac{35}{6}\right), \left(9, -\frac{23}{6}\right), \left(\frac{46}{3}, -\frac{35}{6}\right),$ and $\left(18, -\frac{23}{2}\right)$, respectively. We'll have to throw in the projection of the centroid of $\bigtriangleup JAB$ to the line of reflection to apply shoelace, and that point is $\left( \frac{8}{3}, -\frac{23}{2}\right)$

 Finally, applying Shoelace, we get: $\left|\left(\frac{8}{3}\cdot (-\frac{23}{6})+9\cdot (-\frac{35}{6})+\frac{46}{3}\cdot (-\frac{23}{2})+18\cdot (\frac{-23}{2})+\frac{8}{3}\cdot (-\frac{35}{6})\right) - \left((-\frac{35}{6}\cdot 9) +\\(-\frac{23}{6}\cdot \frac{46}{3})+ (-\frac{35}{6}\cdot 18)+(\frac{-23}{2}\cdot \frac{8}{3})+(-\frac{23}{2}\cdot \frac{8}{3})\right)\right|$$=\left|\left(-\frac{92}{9}-\frac{105}{2}-\frac{529}{3}-207-\frac{140}{9}\right)-\left(-\frac{105}{2}-\frac{529}{9}-105-\frac{92}{3}-\frac{92}{3}\right)\right|$$=\left|-\frac{232}{9}-\frac{1373}{6}-207+\frac{529}{9}+\frac{184}{3}+105+\frac{105}{2}\right|$$=\left|\frac{297}{9}-\frac{690}{6}-102\right|=\left| 33-115-102\right|=\left|-184\right|=\boxed{184}$

 Solution by ktong

 note: for a slightly simpler calculation, notice that the heptagon can be divided into two trapezoids of equal area and a small triangle.

 



\textbf{Solution 2 (Homothety)}

Draw the heptagon whose vertices are the midpoints of octagon $ABCDEFGH$ except $J$.We have a homothety since:

 1. $J$ passes through corresponding vertices of the two heptagons.

 2. By centroid properties, our ratio between the sidelengths is $\frac{2}{3},$ and their area ratio is hence $\frac{4}{9}.$ 

 Compute the area of the large heptagon by dividing into two congruent trapezoids and a triangle. The area of each trapezoid is 
 \[= \frac{1}{2}(17+23)\cdot \frac{19}{2}=190.\] The area of each triangle is 
 \[=  \frac{1}{2}\cdot 17\cdot 4=34.\]

 Hence, the area of the large heptagon is 
 \[2\cdot 190+34=414.\] Then, from our homothety, the area of the required heptagon is 
 \[\frac{4}{9}\cdot 414=\boxed{184}.\]~novus677

 \textbf{Supplement} Note that we use $17.$ A proof of this is as follows:

 
                \begin{itemize}
                    \item First, note that the shaded triangles are 6-8-10 triangles. Because the homothety is defined on the midpoints of the sides of the octagon, the smaller triangle created by that midpoint and the surrounding rectangle has sides 3-4-5. Thus, the sidelength of the octagon (17) is just 23-2*3=17.
                \end{itemize}


                



\textbf{Solution 3 (Less bashy finish than shoelace)}

Instead of bashing shoelace, we can find a clever way to calculate the area of the heptagon without using homothety. By connecting the centroids of $\bigtriangleup JAB, \bigtriangleup JCD$, $\bigtriangleup JHG, \bigtriangleup JFE$, and $\bigtriangleup JFE, \bigtriangleup JCD$. This will give us three triangles and a rectangle. The area of the rectangle is $\frac{34}{3}*\frac{38}{3}$, the top and bottom triangles each give $\frac{38}{3}*2*\frac{1}{2}$, and the triangle on the right yields $\frac{34}{3}*\frac{8}{3}*\frac{1}{2}$ by using $\frac{bh}{2}$. Summing them up, we get $\frac{1656}{9}$, which gives us $\boxed{184}$.~boppitybobii

 


	\item (\textit{AIME II 2018}) Find the number of functions $f(x)$ from $\{1, 2, 3, 4, 5\}$ to $\{1, 2, 3, 4, 5\}$ that satisfy $f(f(x)) = f(f(f(x)))$ for all $x$ in $\{1, 2, 3, 4, 5\}$.


 
 \textbf{Answer: }756



\textbf{Solution 1}

Just to visualize solution 1. If we list all possible $(x,f(x))$, from ${1,2,3,4,5}$ to ${1,2,3,4,5}$ in a specific order, we get $5 \cdot 5 = 25$ different $(x,f(x))$ 's.Namely:

 
 \[(1,1) (1,2) (1,3) (1,4) (1,5)\]
 \[(2,1) (2,2) (2,3) (2,4) (2,5)\] 
 \[(3,1) (3,2) (3,3) (3,4) (3,5)\] 
 \[(4,1) (4,2) (4,3) (4,4) (4,5)\]
 \[(5,1) (5,2) (5,3) (5,4) (5,5)\]

 To list them in this specific order makes it a lot easier to solve this problem. We notice, In order to solve this problem at least one pair of $(x,x)$ where $x\in{1,2,3,4,5}$ must exist.In this case I rather "go backwards". First fixing $5$ pairs $(x,x)$, (the diagonal of our table) and map them to the other fitting pairs $(x,f(x))$. You can do this in $\frac{5!}{5!} = 1$ way. Then fixing $4$ pairs $(x,x)$ (The diagonal minus $1$) and map them to the other fitting pairs $(x,f(x))$. You can do this in $4\cdot\frac{5!}{4!} = 20$ ways. Then fixing $3$ pairs $(x,x)$ (The diagonal minus $2$) and map them to the other fitting pairs $(x,f(x))$. You can do this in $\frac{(5\cdot4\cdot3\cdot6\cdot3)}{3!2!} + \frac{(5\cdot4\cdot3\cdot6\cdot1)}{3!} = 150$ ways.Fixing $2$ pairs $(x,x)$ (the diagonal minus $3$) and map them to the other fitting pairs $(x,f(x))$. You can do this in $\frac{(5\cdot4\cdot6\cdot4\cdot2)}{2!3!} + \frac{(5\cdot4\cdot6\cdot4\cdot2)}{2!2!} + \frac{(5\cdot4\cdot6\cdot2\cdot1)}{2!2!} = 380$ ways.Lastly, fixing $1$ pair $(x,x)$ (the diagonal minus $4$) and map them to the other fitting pairs $(x,f(x))$. You can do this in $\frac{5!}{4!} + 4\cdot\frac{5!}{3!} + 5! = 205$ ways.

 So $1 + 20 + 150 + 380 + 205 = \framebox{756}$

 



\textbf{Solution 2}

We perform casework on the number of fixed points (the number of points where $f(x) = x$). There must be at least one fixed point, given $f(f(x))$ has some value and is a fixed point. To better visualize this, use the grid from Solution 1. 

 \textbf{Case 1:} 5 fixed points

 
            \begin{enumerate}[label=(\roman*)]
                \item - Obviously, there must be $1$ way to do so.
            \end{enumerate}
            \textbf{Case 2:} 4 fixed points

 
            \begin{enumerate}[label=(\roman*)]
                \item - $\binom 54$ ways to choose the $4$ fixed points.  For the sake of conversation, let them be $(1, 1), (2, 2), (3, 3), (4, 4)$.\item - The last point must be $(5, 1), (5, 2), (5, 3),$ or $(5, 4)$.\item - There are $\binom 54 \cdot 4 = 20$ solutions for this case.
            \end{enumerate}
            \textbf{Case 3:} 3 fixed points

 
            \begin{enumerate}[label=(\roman*)]
                \item - $\binom 53$ ways to choose the $3$ fixed points.  For the sake of conversation, let them be $(1, 1), (2, 2), (3, 3)$.<dd>- \textbf{Subcase 3.1:} None of $4$ or $5$ map to each other
            \begin{enumerate}[label=(\roman*)]
                \item - The points must be $(4, 1), (4, 2), (4, 3)$ and $(5, 1), (5, 2), (5, 3)$, giving $3 \cdot 3 = 9$ solutions.
            \end{enumerate}
            </dd><dd>- \textbf{Subcase 3.2:} $4$ maps to $5$ or $5$ maps to $4$
            \begin{enumerate}[label=(\roman*)]
                \item - The points must be $(4, 5)$ and $(5, 1), (5, 2), (5, 3)$ or $(5, 4)$ and $(4, 1), (4, 2), (4, 3)$, giving $3+3 = 6$ solutions.
            \end{enumerate}
            </dd>\item - There are $\binom 53 \cdot (9+6) = 150$ solutions for this case.
            \end{enumerate}
            \textbf{Case 4:} 2 fixed points

 
            \begin{enumerate}[label=(\roman*)]
                \item - $\binom 52$ ways to choose the $2$ fixed points.  For the sake of conversation, let them be $(1, 1), (2, 2)$.<dd>- \textbf{Subcase 4.1:} None of $3$, $4$, or $5$ map to each other
            \begin{enumerate}[label=(\roman*)]
                \item - There are $2 \cdot 2 \cdot 2 = 8$ solutions for this case, by similar logic to \textbf{Subcase 3.1}.
            \end{enumerate}
            </dd><dd>- \textbf{Subcase 4.2:} exactly one of $3, 4, 5$ maps to another.
            \begin{enumerate}[label=(\roman*)]
                \item - Example: $(3, 4), (4, 1), (5, 2)$\item - $\binom 32 \cdot 2! = 6$ ways to choose the 2 which map to each other, and $2\cdot 2$ ways to choose which ones of $1$ and $2$ they map to, giving $24$ solutions for this case.
            \end{enumerate}
            </dd><dd>- \textbf{Subcase 4.3:} two of $3, 4, 5$ map to the third
            \begin{enumerate}[label=(\roman*)]
                \item - Example: $(3, 5), (4, 5), (5, 2)$\item - $\binom 31$ ways to choose which point is being mapped to, and $2$ ways to choose which one of $1$ and $2$ it maps to, giving $6$ solutions for this case.
            \end{enumerate}
            </dd>\item - There are $\binom 52 \cdot (8+24+6) = 380$ solutions.
            \end{enumerate}
            \textbf{Case 5:} 1 fixed point

 
            \begin{enumerate}[label=(\roman*)]
                \item - $\binom 51$ ways to choose the fixed point.  For the sake of conversation, let it be $(1, 1)$.<dd>- \textbf{Subcase 5.1:} None of $2, 3, 4, 5$ map to each other
            \begin{enumerate}[label=(\roman*)]
                \item - Obviously, there is $1^4 = 1$ solution to this; all of them map to $1$.
            \end{enumerate}
            </dd><dd>- \textbf{Subcase 5.2:} One maps to another, and the other two map to $1$.
            \begin{enumerate}[label=(\roman*)]
                \item - Example: $(2, 3), (3, 1), (4, 1), (5, 1)$\item - There are $\binom 42 \cdot 2!$ ways to choose which two map to each other, and since each must map to $1$, this gives $12$.
            \end{enumerate}
            </dd><dd>- \textbf{Subcase 5.3:} One maps to another, and of the other two, one maps to the other as well.
            \begin{enumerate}[label=(\roman*)]
                \item - Example: $(2, 3), (3, 1), (5, 4), (4, 1)$\item - There are $\binom 42 \cdot 2! \cdot 2! / 2!$ ways to choose which ones map to another.  This also gives $12$.
            \end{enumerate}
            </dd><dd>- \textbf{Subcase 5.4:} 2 map to a third, and the fourth maps to $1$.
            \begin{enumerate}[label=(\roman*)]
                \item - Example: $(4, 2), (5, 2), (2, 1), (3, 1)$\item - There are $\binom 42 \cdot \binom 21 = 12$ ways again.
            \end{enumerate}
            </dd><dd>- \textbf{Subcase 5.5:} 3 map to the fourth.
            \begin{enumerate}[label=(\roman*)]
                \item - Example: $(2, 4), (3, 4), (5, 4), (4, 1)$\item - There are $\binom 41$ ways to choose which one is being mapped to, giving $4$ solutions.
            \end{enumerate}
            </dd>\item - There are $\binom 51 \cdot (1+12+12+12+4) = 205$ solutions.
            \end{enumerate}
            Therefore, the answer is $1+20+150+380+205 = \boxed{756}$

 



\textbf{Solution 3}

We can do some caseworks about the special points of functions $f$ for $x\in\{1,2,3,4,5\}$. Let $x$, $y$ and $z$ be three different elements in set $\{1,2,3,4,5\}$. There must be elements such like $k$ in set $\{1,2,3,4,5\}$ satisfies $f(k)=k$, and we call the points such like $(k,k)$ on functions $f$ are "Good Points" (Actually its academic name is "fixed-points"). The only thing we need to consider is the "steps" to get  "Good Points". Notice that the "steps" must less than $3$ because the highest iterations of function $f$ is $3$. Now we can classify $3$ cases of “Good points” of $f$.

 $\textbf{Case 1:}$ One "step" to "Good Points": Assume that $f(x)=x$, then we get $f(f(x))=f(x)=x$, and $f(f(f(x)))=f(f(x))=f(x)=x$, so $f(f(f(x)))=f(f(x))$.

 $\textbf{Case 2:}$ Two "steps" to "Good Points": Assume that $f(x)=y$ and $f(y)=y$, then we get $f(f(x))=f(y)=y$, and $f(f(f(x)))=f(f(y))=f(y)=y$, so $f(f(f(x)))=f(f(x))$.

 $\textbf{Case 3:}$ Three "steps" to "Good Points": Assume that $f(x)=y$, $f(y)=z$ and $f(z)=z$, then we get $f(f(x))=f(y)=z$, and $f(f(f(x)))=f(f(y))=f(z)=z$, so $f(f(f(x)))=f(f(x))$.

 Divide set $\{1,2,3,4,5\}$ into three parts which satisfy these three cases, respectively. Let the first part has $a$ elements, the second part has $b$ elements and the third part has $c$ elements, it is easy to see that $a+b+c=5$. First, there are $\binom{5}{a}$ ways to select $x$ for Case 1. Second, we have $\binom{5-a}{b}$ ways to select $x$ for Case 2. After that we map all elements that satisfy Case 2 to Case 1, and the total number of ways of this operation is $a^b$. Finally, we map all the elements that satisfy Case 3 to Case 2, and the total number of ways of this operation is $b^c$. As a result, the number of such functions $f$ can be represented in an algebraic expression contains $a$, $b$ and $c$: $\boxed{\binom{5}{a}\cdot \binom{5-a}{b}\cdot a^b\cdot b^c}$

 Now it's time to consider about the different values of $a$, $b$ and $c$ and the total number of functions $f$ satisfy these values of $a$, $b$ and $c$:

 For $a=5$, $b=0$ and $c=0$, the number of $f$s is $\binom{5}{5}=1$

 For $a=4$, $b=1$ and $c=0$, the number of $f$s is $\binom{5}{4}\cdot \binom{1}{1}\cdot 4^1\cdot 1^0=20$

 For $a=3$, $b=1$ and $c=1$, the number of $f$s is $\binom{5}{3}\cdot \binom{2}{1}\cdot 3^1\cdot 1^1=60$

 For $a=3$, $b=2$ and $c=0$, the number of $f$s is $\binom{5}{3}\cdot \binom{2}{2}\cdot 3^2\cdot 2^0=90$

 For $a=2$, $b=1$ and $c=2$, the number of $f$s is $\binom{5}{2}\cdot \binom{3}{1}\cdot 2^1\cdot 1^2=60$

 For $a=2$, $b=2$ and $c=1$, the number of $f$s is $\binom{5}{2}\cdot \binom{3}{2}\cdot 2^2\cdot 2^1=240$

 For $a=2$, $b=3$ and $c=0$, the number of $f$s is $\binom{5}{2}\cdot \binom{3}{3}\cdot 2^3\cdot 3^0=80$

 For $a=1$, $b=1$ and $c=3$, the number of $f$s is $\binom{5}{1}\cdot \binom{4}{1}\cdot 1^1\cdot 1^3=20$

 For $a=1$, $b=2$ and $c=2$, the number of $f$s is $\binom{5}{1}\cdot \binom{4}{2}\cdot 1^2\cdot 2^2=120$

 For $a=1$, $b=3$ and $c=1$, the number of $f$s is $\binom{5}{1}\cdot \binom{4}{3}\cdot 1^3\cdot 3^1=60$

 For $a=1$, $b=4$ and $c=0$, the number of $f$s is $\binom{5}{1}\cdot \binom{4}{4}\cdot 1^4\cdot 4^0=5$

 Finally, we get the total number of function $f$, the number is $1+20+60+90+60+240+80+20+120+60+5=\boxed{756}$

 \textbf{Note (fun fact)} This exact problem showed up earlier on the 2011 Stanford Math Tournament, Advanced Topics Test.This problem also showed up on the 2010 Mock AIME 2 here: \href{https://artofproblemsolving.com/wiki/index.php/Mock\_AIME\_2\_2010\_Problems}{https://artofproblemsolving.com/wiki/index.php/Mock\_AIME\_2\_2010\_Problems}

 


	\item (\textit{AIME II 2018}) Find the number of permutations of $1, 2, 3, 4, 5, 6$ such that for each $k$ with $1$ $\leq$ $k$ $\leq$ $5$, at least one of the first $k$ terms of the permutation is greater than $k$.


 
 \textbf{Answer: }461



\textbf{Solution 1}

If the first number is $6$, then there are no restrictions. There are $5!$, or $120$ ways to place the other $5$ numbers.

 
If the first number is $5$, $6$ can go in four places, and there are $4!$ ways to place the other $4$ numbers. $4 \cdot 4! = 96$ ways.

 
If the first number is $4$, ....

 4 6 \_ \_ \_ \_ $\implies$ 24 ways

 4 \_ 6 \_ \_ \_ $\implies$ 24 ways

 4 \_ \_ 6 \_ \_ $\implies$ 24 ways

 4 \_ \_ \_ 6 \_ $\implies$ 5 must go between $4$ and $6$, so there are $3 \cdot 3! = 18$ ways.

 $24 + 24 + 24 + 18 = 90$ ways if 4 is first.

 
If the first number is $3$, ....

 3 6 \_ \_ \_ \_ $\implies$ 24 ways

 3 \_ 6 \_ \_ \_ $\implies$ 24 ways

 3 1 \_ 6 \_ \_ $\implies$ 4 ways 

 3 2 \_ 6 \_ \_ $\implies$ 4 ways

 3 4 \_ 6 \_ \_ $\implies$ 6 ways

 3 5 \_ 6 \_ \_ $\implies$ 6 ways

 3 5 \_ \_ 6 \_ $\implies$ 6 ways

 3 \_ 5 \_ 6 \_ $\implies$ 6 ways

 3 \_ \_ 5 6 \_ $\implies$ 4 ways

 $24 + 24 + 4 + 4 + 6 + 6 + 6 + 6 + 4 = 84$ ways

 
If the first number is $2$, ....

 2 6 \_ \_ \_ \_ $\implies$ 24 ways

 2 \_ 6 \_ \_ \_ $\implies$ 18 ways 

 2 3 \_ 6 \_ \_ $\implies$ 4 ways

 2 4 \_ 6 \_ \_ $\implies$ 6 ways

 2 5 \_ 6 \_ \_ $\implies$ 6 ways

 2 5 \_ \_ 6 \_ $\implies$ 6 ways

 2 \_ 5 \_ 6 \_ $\implies$ 4 ways

 2 4 \_ 5 6 \_ $\implies$ 2 ways

 2 3 4 5 6 1 $\implies$ 1 way

 
$24 + 18 + 4 + 6 + 6 + 6 + 4 + 2 + 1 = 71$ ways

 
Grand Total : $120 + 96 + 90 + 84 + 71 = \boxed{461}$

 



\textbf{Solution 2}

If $6$ is the first number, then there are no restrictions. There are $5!$, or $120$ ways to place the other $5$ numbers.

 
If $6$ is the second number, then the first number can be $2, 3, 4,$ or $5$, and there are $4!$ ways to place the other $4$ numbers. $4 \cdot 4! = 96$ ways.

 
If $6$ is the third number, then we cannot have the following:

 1 \_ 6 \_ \_ \_ $\implies$ 24 ways

 2 1 6 \_ \_ \_ $\implies$ 6 ways

 $120 - 24 - 6 = 90$ ways

 If $6$ is the fourth number, then we cannot have the following:

 1 \_ \_ 6 \_ \_ $\implies$ 24 ways

 2 1 \_ 6 \_ \_ $\implies$ 6 ways

 2 3 1 6 \_ \_ $\implies$ 2 ways

 3 1 2 6 \_ \_ $\implies$ 2 ways

 3 2 1 6 \_ \_ $\implies$ 2 ways

 $120 - 24 - 6 - 2 - 2 - 2 = 84$ ways

 If $6$ is the fifth number, then we cannot have the following:

 \_ \_ \_ \_ 6 5 $\implies$ 24 ways

 1 5 \_ \_ 6 \_ $\implies$ 6 ways

 1 \_ 5 \_ 6 \_ $\implies$ 6 ways

 2 1 5 \_ 6 \_ $\implies$ 2 ways

 1 \_ \_ 5 6 \_ $\implies$ 6 ways

 2 1 \_ 5 6 \_ $\implies$ 2 ways

 2 3 1 5 6 4, 3 1 2 5 6 4, 3 2 1 5 6 4 $\implies$ 3 ways

 $120 - 24 - 6 - 6 - 2 - 6 - 2 - 3 = 71$ ways

 Grand Total : $120 + 96 + 90 + 84 + 71 = \boxed{461}$

 



\textbf{Solution 3 (Recursion)}

Note the condition in the problem is equivalent to the following condition: for each $k$ with $1 \le k \le 5$, the first $k$ terms is not a permutation $(1, 2, \ldots, k)$ (since it would mean there must be some integer $x$ in the first $k$ terms such that $x \not \in \{1, 2, \ldots, k\} \implies x > k$). Then, let $a_n$ denote the number of permutations of $(1, 2, \ldots, n)$ satisfying the condition in the problem. We use complementary counting to find $a_n$. Notice that in order to not satisfy the condition in the problem, there are $n-1$ cases: the first $1 \le k \le n-1$ (we don't include $k = n$ since the condition in the problem only holds up to $n-1$) terms are a permutation of $(1, 2, \ldots, k)$, and for all $k+1 \le i \le n-1$, the condition in the problem still holds. Then, for each of these cases, there are $k!$ ways to arrange the first $k$ terms, and then $a_{n-k}$ ways to arrange the $k + 1$ to $n$ terms (assume by contradiction that terms from $k+1$ to $i$ is a permutation of $(k+1, k+2, \ldots, i)$. Then, since the first $k$ terms are already a permutation of $(1, 2, \ldots, k)$, the first $i$ terms form a permutation of $(1, 2, \ldots, i)$. This contradicts our assumption that for all $k+1 \le i \le n-1$, the condition still holds. Thus, we can only include the $a_{n-k}$ permutations of these terms). Now, summing the cases up and subtracting from $n!$, we have: $a_n = n! - \sum_{k=1}^{n-1} a_{n-k} k!$. From this recursion, we derive that $a_1 = 1$, $a_2 = 1$, $a_3 = 3$, $a_4 = 13$, $a_5 = 71$, and finally $a_6 = \boxed{461}$.

 



\textbf{Solution 4 (PIE)}

Let $A_i$ be the set of permutations such that there is no number greater than $i$ in the first $i$ places. Note that $\bigcap^{k}_{i=0}{A_{b_i}}=\prod^k_{i=1}{(b_i-b_{i-1})!}$ for all $1\le b_0 < b_1\cdots < b_k \le 5$ and that the set of restricted permutations is $A_1 \cup A_2 \cup A_3 \cup A_4 \cup A_5$.

 We will compute the cardinality of this set with PIE.
 \begin{align*} &|A_1| + |A_2| + |A_3| + |A_4| + |A_5|\\ = &120 + 48 + 36 + 48 + 120 = 372\\ \\ &|A_1 \cap A_2| + |A_1 \cap A_3| + |A_1 \cap A_4| + |A_1 \cap A_5| + |A_2 \cap A_3|\\ + &|A_2 \cap A_4| + |A_2 \cap A_5| + |A_3 \cap A_4| + |A_3 \cap A_5| + |A_4 \cap A_5|\\=&24+12+12+24+12+8+12+12+12+24=152\\ \\ &|A_1 \cap A_2 \cap A_3| + |A_1 \cap A_2 \cap A_4| + |A_1 \cap A_2 \cap A_5| + |A_1 \cap A_3 \cap A_4| + |A_1 \cap A_3 \cap A_5|\\ +& |A_1 \cap A_4 \cap A_5| + |A_2 \cap A_3 \cap A_4| + |A_2 \cap A_3 \cap A_5| +  |A_2 \cap A_4 \cap A_5| + |A_3 \cap A_4 \cap A_5|\\=&6 + 4 + 6 + 4 + 4 + 6 + 4 + 4 + 4 + 6 = 48\\ \\ &|A_1 \cap A_2 \cap A_3 \cap A_4| + |A_1 \cap A_2 \cap A_3 \cap A_5| + |A_1 \cap A_2 \cap A_4 \cap A_5| + |A_1 \cap A_3 \cap A_4 \cap A_5| + |A_2 \cap A_3 \cap A_4 \cap A_5|\\=&2 + 2 + 2 + 2 + 2 = 10\\ \\ &|A_1 \cap A_2 \cap A_3 \cap A_4 \cap A_5| = 1 \end{align*}
To finish, $720 - 372 + 152 - 48 + 10 - 1 = \boxed{461}$

 



\textbf{Solution 5 (Recursion)}

Define the function $f(p,q)$ as the amount of permutations with maximum digit $q$ and string length $p$ that satisfy the condition within bounds. For example, $f(4,5)$ would be the amount of ways to make a string with length $4$ with the highest digit being $5$. We wish to obtain $f(6,6)=f(5,6)$. 

 To generate recursion, consider how we would get to $f(p,q)$ from $f(p-1,a)$ for all $a$ such that $p\le{a}\le6$. We could either jump from the old maximum $a$ to the new $q$ by concatenating the old string and the new digit $q$, or one could retain the maximum, in which case $a=q$. To retain the maximum, one would have to pick a new available digit not exceeding $q$. 

 In the first case, there is only one way to pick the new digit, namely picking $q$. For the second case, there are $q-p+1$ digits left to choose, because there are $q$ digits between 1 and $q$ total and there are $p-1$ digits already chosen below or equal to $q$. Thus, $f(p,q)=[\sum^{q-1}_{n=p}f(p-1,n)] + (q-p+1)f(p-1,q)$. Now that we have the recursive function, we can start evaluating the values of $f(p,q)$ until we get to $f(6,6)=f(5,6)$. 

 
 \[f(2,3)=3, f(2,4)=5, f(2,5)=7, f(2,6)=9\]
 \[f(3,4)=13, f(3,5)=29, f(3,6)=51\]
 \[f(4,5)=71, f(4,6)=195\]
 \[f(5,6)=461\]Our requested answer is thus $\boxed{461}$~sigma

 



\textbf{Solution 6 (Complementary)}

We can also solve this problem by counting the number of permutations that do NOT satisfy the given conditions; namely, these permutations must satisfy the condition that none of the first $k$ terms is greater than $k$ for $1\leq$ $k$ $\leq$ $5$. We can further simplify this method by approaching it through casework on the first $k$ terms.

 Case 1: None of the first one terms is greater than 1

 The first term must obviously be one. Since there are five spaces left for numbers, there are a total of $5!=120$ permutations for this case.

 Case 2: None of the first two terms is greater than 2

 The first two terms must be 1 and 2 in some order. However, we already counted all cases starting with 1 in the previous case, so all of the permutations in this case must begin with $21\cdots$. Since there are four spaces left, there are a total of $4!=24$ permutations for this case.

 Case 3: None of the first three terms is greater than 3

 The first three terms must be 1, 2, and 3 in some order. However, the cases beginning with 1\_\_ and 21\_ have already been accounted for. There are now $3!-3 = 3$ ways to order the first three numbers of the permutation, and $3!$ ways to order the last three numbers, for a total of $3\times6 = 18$ permutations.

 Case 4: None of the first four terms is greater than 4

 You can see where the pattern is going - the first four terms must be 1, 2, 3, and 4 in some order. All cases starting with 1 (there are $3!=6$), the cases starting with 21 (there are $2!=2$), and the 3 cases from case 3 (there are $3\times 1! = 3$) have been accounted for, so there are a total of $(4!-6-2-3)2!=26$ permutations for this case.

 Case 5: None of the first five terms is greater than 5

 This is perhaps the hardest case to work with, simply because there are so many subcases, so keeping track is crucial here. Obviously, the first five terms must be 1, 2, 3, 4, and 5, meaning there are 120 ways to order them. Now, we count the permutations we have already counted in previous cases. $4!$ start with 1, $3!$ start with 2, $3\times2!=6$ start with 3, and $13\times1!=13$ start with 4. Subtracting, we get a total of $120-24-6-6-13=71$ permutations.

 Now, we subtract the total number of permutations from our cases from the total number of permutations, which is $6!$ :$720 - 120 - 24 - 18 - 26 - 71 = \boxed{461}$.

 



\textbf{Solution 7(Similar to Solutions 1 and 2 but with more depth)}

We will perform cases based on the position of $6$. Notice that if $6$ is the first number, the positions of the rest of the numbers absolutely don't matter, so we have $5!$ cases here.

 Now, if $6$ is the second number, all we need is the first number to be greater than $1$. Namely, the possibilities for the first number are just $2$ to $5$, for a total of $4$ numbers. The rest of the number positions don't matter, so we have $4!$ ways to arrange the rest of the numbers. Hence, for this case, we have $4 \cdot 4!$ cases. 

 Now, if $6$ is the third number, we can break it up into two cases. First, the first number is simply $2$. Then, as "at least one of the first $2$ terms is greater than $2$" we just need the second term to be greater than $2$, namely the possibilities are $3$ to $5$, a total of $3$ numbers. The arrangement of the remaining numbers doesn't matter, and hence that is just $3!$. The first part of this case yields $3 \cdot 3!$. Now, the second part is that the first term is any number from $3$ to $5$. Then, the rest of the numbers don't matter. Hence, we have $4!$ ways to arrange the rest of the numbers. So we have a total of $3! \cdot 3 + 4! \cdot 3$ ways for this case. 

 Now, if $6$ is the fourth number, things get a little bit more difficult. We will break it up into $4$ cases. First, we have $2$ as the first term. Now, if the second term is $3$, all we need for this to succeed is for the third term to be at least $4$, so only two possibilities $4$ and $5$. The rest of the number arrangements don't matter, so we have $2!$ ways for that. So the first part of this case is $2 \cdot 2!$. Now, if the second term is at least $4$, then we are done. Hence, we have our two possibilities $4$ and $5$ for the second term, and the remaining number arrangements don't matter. Hence, we have $2 \cdot 3!$. So if the first term if $2$, we have $2 \cdot 2! + 2 \cdot 3!$ cases.Now, if the first term is $3$, all we need to satisfy is that "at least one of the first $3$ terms is greater than $3$." Hence, we can break this up also into cases. First, we'll consider two blocks: one block of numbers before the $6$ and one block of numbers after the $6$. If $4$ is in the first block before $6$ and $5$ is in the second block, we have $2!$ ways to arrange the rest of the numbers. There are $2$ ways to place the $4$ in the first block and $2$ ways to place the $5$ in the second block for a total of $2 \cdot 2 = 4$. We can also switch the positions of $4$ and $5$, namely, placing $5$ in the first block and $4$ in the second block, as this won't affect the requirements. So, we have a total of $2! \cdot 2 \cdot 4$ ways for the first part. Now, the second part is when both $4$ and $5$ are in the first block. There are $2$ ways for both $4$ and $5$ to be in the first block and then $2!$ ways to arrange the rest of the numbers. So we have $2! \cdot 2$ for this case. Hence, if $3$ is the first term, we have $2! \cdot 2 \cdot 4 + 2! \cdot 2$ cases.Now, if $4$ is the first term, everything is satisfied, and we are left with the task of finding how many ways there are to arrange the rest of the numbers. This is simply $4!$. The same occurs if $5$ is the first term with $4!$ ways. 

 After all the casework, if $6$ is the fourth number, there is a total of $2 \cdot 2! + 2 \cdot 3! + 2! \cdot 2 \cdot 4 + 2! \cdot 2 + 4! + 4!$. 

 Finally, we move on to the last case of when $6$ is the fifth number. We can do a similar casework analysis as we did to find the previous case. We will break up the cases based on which number is placed in the first term. 

 First, consider if $2$ is the first term. Now, if the second term is $3$, then we need $4$ and $5$ to be both placed in the first block. Otherwise, if $4$ and $5$ are placed in different blocks or they're both in the second block, the conditions won't be satisfied. If the third term is $4$, then $5$ has to be in the fourth term, and we have only $1$ way for this to happen. Now, if the third term is $5$, everything is satisfied, and we have $2!$ ways to arrange them. So if $3$ is the second term, we have $1 + 2!$ ways to arrange them. Now, if the second term is $4$, $5$ has to be placed in the first block, which has only $2$ remaining positions. After placing the $5$, we have $2$ remaining spots overall, so we have $2! \cdot 2$ ways for this to happen. Finally, if $5$ is the second term, everything is satisfied, so we have $3!$ ways to arrange the rest of the numbers.If $2$ is the first term, we have $3! + 2! \cdot 2 + 1 + 2!$ ways. 

 Now, let's consider $3$ being the first term. We will do this case a little differently. We will consider the position of the number $5$. Let's say $5$ is the second term. Then everything is satisfied, and we have $3!$ ways for this to happen. If $5$ is the third term, everything is once again satisfied, and we have $3!$ ways. If $5$ is the fourth term, we need the number $4$ to be either in the second or third positions, with, after the $4$ is placed, leaving $2$ remaining positions with $2!$ ways to arrange them. Hence, we have $2 \cdot 2!$. If $3$ is the first term, we have $3! + 3! + 2 \cdot 2!$ ways.

 Now, if $4$ is the first term, all we need is to satisfy "at least one of the first $4$ terms is greater than $4$." Therefore, $5$ must be placed in the first block where there are $3$ positions. The order of the rest of the numbers doesn't matter. We have $3! \cdot 3$ ways for this to happen.

 Finally, if $5$ is the first term, everything is satisfied, and we have $4!$ ways for this to happen.

 Finally, the grand total is

 $5! + 4! \cdot 4 + 3! \cdot 3 + 4! \cdot 3 + 2 \cdot 2! + 2 \cdot 3! + 2! \cdot 2 \cdot 4 + 2! \cdot 2 + 4! + 4! + 3! + 2 \cdot 2! + 1 + 2! + 3! + 3! + 2 \cdot 2! + 3 \cdot 3! + 4! = \boxed{461}$.

 


	\item (\textit{AIME II 2018}) Let $ABCD$ be a convex quadrilateral with $AB = CD = 10$, $BC = 14$, and $AD = 2\sqrt{65}$. Assume that the diagonals of $ABCD$ intersect at point $P$, and that the sum of the areas of triangles $APB$ and $CPD$ equals the sum of the areas of triangles $BPC$ and $APD$. Find the area of quadrilateral $ABCD$.


 
 \textbf{Answer: }112



\textbf{Solution 1}

Let $AP=x$ and let $PC=\rho x$. Let $[ABP]=\Delta$ and let $[ADP]=\Lambda$. We easily get $[PBC]=\rho \Delta$ and $[PCD]=\rho\Lambda$. 

 We are given that $[ABP] +[PCD] = [PBC]+[ADP]$, which we can now write as 
 \[\Delta + \rho\Lambda = \rho\Delta + \Lambda \qquad \Longrightarrow \qquad \Delta -\Lambda = \rho (\Delta -\Lambda).\] Either $\Delta = \Lambda$ or $\rho=1$. The former would imply that $ABCD$ is a parallelogram, which it isn't; therefore we conclude $\rho=1$ and $P$ is the midpoint of $AC$. Let $\angle BAD = \theta$ and $\angle BCD = \phi$. Then $[ABCD]=2\cdot [BCD]=140\sin\phi$. On one hand, since $[ABD]=[BCD]$, we have 
 \begin{align}\sqrt{65}\sin\theta = 7\sin\phi \quad \implies \quad 16+49\cos^2\phi = 65\cos^2\theta\end{align}whereas, on the other hand, using cosine formula to get the length of $BD$, we get 
 \[10^2+4\cdot 65 - 40\sqrt{65}\cos\theta = 10^2+14^2-280\cos\phi\]
 \begin{align}\tag{2}\implies \qquad  65\cos^2\theta = \left(7\cos\phi+ \frac{8}{5}\right)^2\end{align}Eliminating $\cos\theta$ in the above two equations and solving for $\cos\phi$ we get
 \[\cos\phi = \frac{3}{5}\qquad \implies \qquad \sin\phi = \frac{4}{5}\]which finally yields $[ABCD]=2\cdot [BCD] = 140\sin\phi = 112$.

 



\textbf{Solution 2}

For reference, $2\sqrt{65} \approx 16$, so $\overline{AD}$ is the longest of the four sides of $ABCD$. Let $h_1$ be the length of the altitude from $B$ to $\overline{AC}$, and let $h_2$ be the length of the altitude from $D$ to $\overline{AC}$. Then, the triangle area equation becomes

 
 \[\frac{h_1}{2}AP + \frac{h_2}{2}CP = \frac{h_1}{2}CP + \frac{h_2}{2}AP \rightarrow \left(h_1 - h_2\right)AP = \left(h_1 - h_2\right)CP \rightarrow AP = CP\]

 What an important finding! Note that the opposite sides $\overline{AB}$ and $\overline{CD}$ have equal length, and note that diagonal $\overline{DB}$ bisects diagonal $\overline{AC}$. This is very similar to what happens if $ABCD$ were a parallelogram with $AB = CD = 10$, so let's extend $\overline{DB}$ to point $E$, such that $AECD$ is a parallelogram. In other words, 
 \[AE = CD = 10\] and 
 \[EC = DA = 2\sqrt{65}\] Now, let's examine $\triangle ABE$. Since $AB = AE = 10$, the triangle is isosceles, and $\angle ABE \cong \angle AEB$. Note that in parallelogram $AECD$, $\angle AED$ and $\angle CDE$ are congruent, so $\angle ABE \cong \angle CDE$ and thus 
 \[\text{m}\angle ABD = 180^\circ - \text{m}\angle CDB\] Define $\alpha := \text{m}\angle CDB$, so $180^\circ - \alpha = \text{m}\angle ABD$. 

 We use the Law of Cosines on $\triangle DAB$ and $\triangle CDB$:

 
 \[\left(2\sqrt{65}\right)^2 = 10^2 + BD^2 - 20BD\cos\left(180^\circ - \alpha\right) = 100 + BD^2 + 20BD\cos\alpha\]

 
 \[14^2 = 10^2 + BD^2 - 20BD\cos\alpha\]

 Subtracting the second equation from the first yields

 
 \[260 - 196 = 40BD\cos\alpha \rightarrow BD\cos\alpha = \frac{8}{5}\]

 This means that dropping an altitude from $B$ to some foot $Q$ on $\overline{CD}$ gives $DQ = \frac{8}{5}$ and therefore $CQ = \frac{42}{5}$. Seeing that $CQ = \frac{3}{5}\cdot BC$, we conclude that $\triangle QCB$ is a 3-4-5 right triangle, so $BQ = \frac{56}{5}$. Then, the area of $\triangle BCD$ is $\frac{1}{2}\cdot 10 \cdot \frac{56}{5} = 56$. Since $AP = CP$, points $A$ and $C$ are equidistant from $\overline{BD}$, so 
 \[\left[\triangle ABD\right] = \left[\triangle CBD\right] = 56\] and hence 
 \[\left[ABCD\right] = 56 + 56 = \boxed{112}\] -kgator

 
Just to be complete -- $h_1$ and $h_2$ can actually be equal. In this case, $AP \neq CP$, but $BP$ must be equal to $DP$. We get the same result. -Mathdummy.

 



\textbf{Solution 3 (Another way to get the middle point)}

So, let the area of $4$ triangles $\triangle {ABP}=S_{1}$,  $\triangle {BCP}=S_{2}$,  $\triangle {CDP}=S_{3}$,  $\triangle {DAP}=S_{4}$. Suppose $S_{1}>S_{3}$ and $S_{2}>S_{4}$, then it is easy to show that 
 \[S_{1}\cdot S_{3}=S_{2}\cdot S_{4}.\] Also, because 
 \[S_{1}+S_{3}=S_{2}+S_{4},\] we will have 
 \[(S_{1}+S_{3})^2=(S_{2}+S_{4})^2.\] So 
 \[(S_{1}+S_{3})^2=S_{1}^2+S_{3}^2+2\cdot S_{1}\cdot S_{3}=(S_{2}+S_{4})^2=S_{2}^2+S_{4}^2+2\cdot S_{2}\cdot S_{4}.\] So 
 \[S_{1}^2+S_{3}^2=S_{2}^2+S_{4}^2.\] So 
 \[S_{1}^2+S_{3}^2-2\cdot S_{1}\cdot S_{3}=S_{2}^2+S_{4}^2-2\cdot S_{2}\cdot S_{4}.\] So 
 \[(S_{1}-S_{3})^2=(S_{2}-S_{4})^2.\] As a result, 
 \[S_{1}-S_{3}=S_{2}-S_{4}.\] Then, we have 
 \[S_{1}+S_{4}=S_{2}+S_{3}.\] Combine the condition 
 \[S_{1}+S_{3}=S_{2}+S_{4},\] we can find out that 
 \[S_{3}=S_{4},\] so $P$ is the midpoint of $\overline {AC}$

 



\textbf{Solution 4 (With yet another way to get the middle point)}

Denote $\angle APB$ by $\alpha$. Then $\sin(\angle APB)=\sin \alpha = \sin(\angle APD)$.Using the formula for the area of a triangle, we get 
 \[\frac{1}{2} (AP\cdot BP+ CP\cdot DP)\sin\alpha=\frac{1}{2}(AP\cdot DP+ CP\cdot BP)\sin\alpha ,\] so 
 \[(AP-CP)(BP-DP)=0\]Hence $AP=CP$ (note that $BP=DP$ makes no difference here).Now, assume that $AP=CP=x$, $BP=y$, and $DP=z$. Using the cosine rule for $\triangle APB$ and $\triangle BPC$, it is clear that 
 \[x^2+y^2-100=2 xy \cdot \cos{APB}=-(2 xy \cdot \cos{(\pi-CPB)})=-(x^2+y^2-196)\] or 
 \begin{align}x^2+y^2=148\end{align}.Likewise, using the cosine rule for triangles $APD$ and $CPD$, 
 \begin{align}\tag{2}x^2+z^2=180\end{align}. It follows that 
 \begin{align}\tag{3}z^2-y^2=32\end{align}. Since $\sin\alpha=\sqrt{1-\cos^2\alpha}$, 
 \[\sqrt{1-\frac{(x^2+y^2-100)^2}{4x^2y^2}}=\sqrt{1-\frac{(x^2+z^2-260)^2}{4x^2z^2}}\] which simplifies to 
 \[\frac{48^2}{y^2}=\frac{80^2}{z^2} \qquad \Longrightarrow \qquad 5y=3z.\] Plugging this back to equations $(1)$, $(2)$, and $(3)$, it can be solved that $x=\sqrt{130},y=3\sqrt{2},z=5\sqrt{2}$. Then, the area of the quadrilateral is 
 \[x(y+z)\sin\alpha=\sqrt{130}\cdot8\sqrt{2}\cdot\frac{14}{\sqrt{260}}=\boxed{112}\]--Solution by MicGu

 



\textbf{Solution 5}

As in all other solutions, we can first find that either $AP=CP$ or $BP=DP$, but it's an AIME problem, we can take $AP=CP$, and assume the other choice will lead to the same result (which is true). 

 From $AP=CP$, we have $[DAP]=[DCP]$, and $[BAP]=[BCP] \implies [ABD] = [CBD]$, therefore,
 \begin{align} \nonumber \frac 12 \cdot AB\cdot AD\sin A &= \frac 12 \cdot BC\cdot CD\sin C \\ \Longrightarrow \hspace{1in} 7\sin C &= \sqrt{65}\sin A  \end{align}By Law of Cosines,
 \begin{align} \nonumber 10^2+14^2-2\cdot 10\cdot 14\cos C &= 10^2+4\cdot 65-2\cdot 10\cdot 2\sqrt{65}\cos A \\ \Longrightarrow \hspace{1in} -\frac 85 - 7\cos C &=  \sqrt{65}\cos A \tag{2} \end{align}Square $(1)$ and $(2)$, and add them, to get
 \[\left(\frac 85\right)^2  + 2\cdot \frac 85 \cdot 7\cos C + 7^2 = 65\]Solve, $\cos C = 3/5  \implies \sin C = 4/5$,
 \[[ABCD] = 2[BCD] = BC\cdot CD\cdot \sin C = 14\cdot 10\cdot \frac 45 = \boxed{112}\]-Mathdummy

 



\textbf{Solution 6}

Either $PA=PC$ or $PD=PB$. Let $PD=PB=s$. Applying Stewart's Theorem on $\triangle ABD$ and $\triangle BCD$, dividing by $2s$ and rearranging, 
 \[\tag{1}CP^2+s^2=148\] 
 \[\tag{2}AP^2+s^2=180\]   Applying Stewart on $\triangle CAB$ and $\triangle CAD$, 
 \[\tag{3} 5CP^2=3AP^2\] Substituting equations 1 and 2 into 3 and rearranging, $s=BP=PD\sqrt{130}, CP=3\sqrt{2}, PA=5\sqrt{2}$ . By Law of Cosines on $\triangle APB$, $\cos(\angle APB)=\frac{4\sqrt{65}}{65}$ so $\sin(\angle APB)=\sin(\angle BPC)=\sin(\angle CPD)=\sin(\angle DPA)=\frac{7\sqrt{65}}{65}$. Using $[\triangle ABC]=\frac{ab\sin(\angle C)}{2}$ to find unknown areas, $[ABCD]=[\triangle APB]+[\triangle BPC]+[\triangle CPD]+[\triangle DPA]=\boxed{112}$.

 -Solution by Garrett

 



\textbf{Solution 7}

Now we prove P is the midpoint of $BD$. Denote the height from $B$ to $AC$ as $h_1$, height from $D$ to $AC$ as $h_2$.According to the problem, $AP* h_1 +CP* h_2 =CP* h_1 +AP* h_2$ implies $h_1 (AP-CP)= h_2 (AP-CP), h_1 = h_2$. Then according to basic congruent triangles we get $BP=DP$Firstly, denote that $CP=a,BP=b,CP=c,AP=d$. Applying Stewart theorem, getting that $100c+196b=(b+c)(bc+a^2); 100b+260c=(b+c)(bc+c^2); 100c+196b=100b+260c, 3b=5c$, denote $b=5x,c=3x$Applying Stewart Theorem, getting $260a+100a=2a(a^2+25x^2); 196a+100a=2a(9x^2+a^2)$ solve for a, getting $a=\sqrt{130},AP=5\sqrt{2}; CP=3\sqrt{2}$Now everything is clear, we can find $cos\angle{BPA}=\frac{4}{\sqrt{65}}$ using LOC, $sin\angle{BPA}=\frac{7\sqrt{65}}{65}$, the whole area is $\sqrt{130}*8\sqrt{2}*\frac{7\sqrt{65}}{65}=\boxed{112}$

 



\textbf{Solution 8 (Simple Geometry)}


                    \begin{figure}[H]
                        \centering
                        \includegraphics[width=0.5\textwidth]{Solution/AIME_II_2018_12_Solution_Figure_1.png}
                    \end{figure}
                    $BP = PD$ as in another solutions.

 Let $D'$ be the reflection of $D$ across $C$.Let points $E, E',$ and $H$ be the foot of perpendiculars on $AC$ from $D,D'$, and $B$ respectively.
 \begin{align*} &\qquad AB = CD = CD', \quad \textrm{and} \quad BH = DE = D'E' \\ \Rightarrow &\qquad \triangle ABH =  \triangle CDE =  \triangle CD'E' \\ \Rightarrow &\qquad \angle BAC = \angle ACD' \\ \Rightarrow &\qquad \triangle ABC =  \triangle AD'C \\ \Rightarrow &\qquad BC = AD'. \end{align*}
The area of quadrilateral $ABCD$ is equal to the area of triangle $ADD'$ with sides $AD' = 14, AD = 2\sqrt{65}, DD' = 2 \cdot 10 = 20$.The semiperimeter is $s = 17 + \sqrt{65},$ the area 
 \[[ADD'] = \sqrt {(17 + \sqrt{65}) (17 - \sqrt{65})(3 + \sqrt{65})(\sqrt{65}-3)} = \sqrt{(289 – 65)\cdot(65-9)} =\sqrt{56 \cdot 4 \cdot 56} = \boxed{112}.\]

 \textbf{vladimir.shelomovskii@gmail.com, vvsss}

 



\textbf{Solution 9 (Mindless Law of Cosines Bash)}

Use your favorite method to get that $P$ is the midpoint of one of the two diagonals (suppose it's the midpoint of $\overline{AC}$). From here, let $x=AP=PC, y=BP, z=PD, a=\cos\theta$ where $\theta$ is the angle that the diagonals make. Then we have a system of four equations:

 
 \begin{align*} x^2+y^2+2xya &= 100 \\ z^2+x^2+2xza &= 100 \\ x^2+y^2-2xya &= 196 \\ x^2+z^2-2xza &= 260 \\ \end{align*}


 From these equations we get that 
 \begin{align*} xya &= -24 \\ xza &= -40 \\ x^2+y^2-48 &= 10 \\ x^2+z^2-80 &= 10 \end{align*}


 From here we can see that $\frac{z}{y}=\frac{5}{3}, z^2-y^2=32,$ so $z=5\sqrt{2}, y=3\sqrt{2}.$ Furthermore, this implies $x=\sqrt{130}$ and $xa=-4\sqrt{2},$ which implies $a=\cos\theta=\frac{4}{\sqrt{65}}.$ Then note that the area of the quadrilateral is 
 \[\frac{1}{2}\sin\theta (xy+xz+xz+xy)=\sin\theta (\sqrt{130}\cdot 3\sqrt{2}+\sqrt{130} \cdot 5\sqrt{2})=7\cdot (3\cdot 2+5\cdot 2)=7(6+10)=7\cdot 16=\boxed{112}.\]

 



\textbf{Solution 10}

Note that $\angle APB = \angle CPD = 180-\angle APD = 180-\angle BPC.$(All angles are in degrees)Since $\sin(\theta)=\sin(180-\theta),$ we can use sine area formula to get the following(after some simplifying steps):
 \begin{align*} BP \times AP + CP \times DP = BP \times PC + AP \times PD. \end{align*}
For convenience, let $AP=a, BP=b, CP=c,DP=d.$ The above equation simplifies to:
 \begin{align*} ab + cd = bc + ad \\ab-ad+cd-bc=0 \\a(b-d)-c(b-d)=0 \\(a-c)(b-d)=0 \end{align*}
From here, we see that $a=c$ or $b=d$. Without loss of generality, let $a=c$. Since triangles $ABP$ and $CDP$ are obviously not congruent, we see that one triangle is obtuse and the other one is acute.(Refer to the diagram) However, if we drop perpendiculars from $B$ to $AC$ and $D$ to $AC$, we do get congruent triangles. If the foot of the perpendicular from $B$ is $M$, and the foot of the perpendicular from $D$ is $N$, then right triangle $BMP$ is congruent to right triangle $DNP$.From here, we see that the altitudes of triangles $ABC$ and $ADC$ to $AC$ are equal. Since they share base $AC$, their areas are equal. We can use Heron's formula. To not have any fractions, let $AC=2x.$
 \begin{align*} \sqrt{(12+x)(12-x)(x+2)(x-2)}=\sqrt{5+\sqrt{65}+x)(5+\sqrt{65}-x)(5-\sqrt{65}+x)(\sqrt{65}-5+x)} \end{align*}
Even though this looks bad at first, it actually isn't too complicated to simplify. Expanding the differences of squares and simplifying completely, we get $x^2=32.$ Plugging this $x$ back into the Heron's formula, we get that the area of $ABC$(or $ADC$) is $56$. Since these triangles have equal area, the area of the quadrilateral is $2 \times 56 = \boxed{112}$, and we are done. $\blacksquare$

 



\textbf{Solution 11}


                    \begin{figure}[H]
                        \centering
                        \includegraphics[width=0.5\textwidth]{Solution/AIME_II_2018_12_Solution_Figure_2.png}
                    \end{figure}
                    Use any method to derive that $P$ is the midpoint of $A$ and $C$. Now, either construct a parallelogram with as shown in the diagram to the right, or use law of sines. We will take the parallelogram route here, but the same thing can be done with law of sines on triangles $\triangle \textnormal{ABP}$ and $\triangle \textnormal{CPD}$. Reflect $D$ across $P$ to get $D'$. Since $CD = AD' = AB = 10$, $\triangle \textnormal{ABD'}$ is isosceles. Thus, $\angle AD'B = \angle ABD'$, and because $ADCD'$ is a parallelogram (since $AP = PC$ and $DP = PD'$), $\angle AD'B = \angle BDC = \angle ABD'$. So, $\angle ABD = 180 - \angle ABD' = 180 - \angle BDC$. Now, apply law of cosines on $\triangle \textnormal{ABD}$ and $\triangle \textnormal{CDB}$. We get:\begin{align}100 + BD^2 - 20BD \cos{\angle ABD} &= 100 + BD^2 - 20 BD \cos {(180 - \angle BDC)} = \\100 + BD^2 + 20 BD \cos{\angle BDC} &= 260 \\&\textnormal{and} \\100 + BD^2 - 20 BD \cos{\angle BDC} &= 196 \\\textnormal{summing }&\textnormal{and simplifying,} \\BD &= 8\sqrt{2}\end{align}Then, applying law of cosines on $\triangle \textnormal{BCD}$ again, we obtain
 \[100 + 196 - 280 \cos{\angle BCD} = BD^2 = 128 \implies \cos{\angle BCD} = \frac{3}{5} \implies \sin{\angle BCD} = \frac{4}{5}\]Since  $AP = PD$, $[ABD] = [BCD] \implies [ABCD] = [ABD] + [BCD] = 2[BCD]$. Thus, $[ABCD] = 2[BCD] = 2 \cdot \frac{1}{2} \cdot 10 \cdot 14 \sin{\angle BCD} = 140 \cdot \frac{4}{5} = \boxed{112}$.

 



\textbf{Solution 12(LoC Bash)}

Let $AP = a, CP = b, BP = c, DP = d$. Also, let $\angle APB = \theta$ and the other angles will follow by Supplementary Angle Theorem and Vertical Angle Theorem. 

 Using Law of Cosines we get

 $a^2 + d^2 + 2ad\cos(\theta) = 260 \textbf{(1)}$

 $b^2 + c^2 + 2bc\cos(\theta) = 196 \textbf{(2)}$

 $a^2 + c^2 - 2ac\cos(\theta) = 100 \textbf{(3)}$

 $b^2 + d^2 - 2bd\cos(\theta) = 100 \textbf{(4)}$

 By sine area formula

 $\frac{1}{2} \sin(\theta) ac + \frac{1}{2} \sin(\theta) bd = \frac{1}{2} \sin(180 - \theta) bc + \frac{1}{2} \sin(180 - \theta) ad \implies ac + bd = bc + ad \implies c(a - b) = d(a - b) \implies \text{Either a = b or c = d}$

 We will assume $a = b$ for this solution.

 Now we subtract equations (1) and (2)

 $a^2 - b^2 + d^2 - c^2 + 2\cos(\theta)(ad - bc) = 64$

 We also subtract equations (3) and (4)

 $a^2 - b^2 + c^2 - d^2 + 2\cos(\theta)(bd - ac) = 0$

 We will add these equations to get

 $2a^2 - 2b^2 + 2\cos(\theta)((d - c)(a + b)) = 64$

 Then factoring an $a - b$ out from each of the factors and dividing by $2$

 $(a + b)(a - b + \cos(\theta)(d - c)) = 32 \textbf{(5)}$

 Now we will subtract equations (1) and (3)

 $d^2 - c^2 + 2\cos(\theta)(ac + ad) = 160$

 We will also subtract equations (2) and (4)

 $c^2 - d^2 + 2\cos(\theta)(bc + bd) = 96$

 We subtract these equations to get

 $2d^2 - 2c^2 + 2\cos(\theta)(ad - bc - bd + ac) = 64$

 Factoring we have

 $(d + c)(d - c + \cos(\theta)(a - b)) = 32 \textbf{(6)}$

 Now remember that $a = b$. So we can plug that in for both equations (5) and (6)

 $(2a)((d - c)\cos(\theta)) = 32$

 $(d + c)(d - c) = 32$

 We know $d - c = \frac{16}{a\cos(\theta)}$ from the first equation and $d - c = \frac{32}{d + c}$. Setting them equal and solving for $\cos(\theta)$ gives

 $\cos(\theta) = \frac{c + d}{2a}$

 Now we will subtract equations (1) and (4)

 $a^2 - b^2 + 2\cos(\theta)(ad + bd) = 160$

 Factoring we have

 $(a + b)(a - b + 2d\cos(\theta)) = 160$

 Substituting $a = b$ we get

 $(a + b)(2d\cos(\theta) = 160 \implies \cos(\theta) = \frac{80}{d(a + b)}$

 Now we will subtract equations (2) and (3)

 $b^2 - a^2 + 2\cos(\theta)(bc + ac) = 96$

 Factoring we get

 $(a + b)(b - a + 2c\cos(\theta)) = 96$

 We can substitute $a = b$ to get

 $(a + b)(2c\cos(\theta)) = 96 \implies \cos(\theta) = \frac{48}{c(a + b)}$

 Now setting these two representations of cosines equal, we have

 $\frac{80}{d} = \frac{48}{c}$ after cancelling the $a + b$'s. 

 Solving, we have $d = \frac{5c}{3}$. 

 Note that $\cos(\theta) = \frac{c + d}{2a}$ from earlier. Plugging in $d = \frac{5c}{3}$, we get

 $\cos(\theta) = \frac{4c}{3a}$

 We have from equation 6

 $(d + c)(d - c + \cos(\theta)(a - b)) = 32$

 We can plug in $a = b$ to get

 $d^2 - c^2 = 32$ following difference of squares.

 Now $d = \frac{5c}{3}$ so we can substitute that in to have

 $\frac{25c^2}{9} - \frac{9c^2}{9} = 32 \implies \frac{16c^2}{9} = 32 \implies c^2 = 18 \implies c = 3\sqrt{2}$ because $c$ has to be positive as it is a side length. 

 Now $d = \frac{5c}{3} = \frac{5 \cdot 3\sqrt{2}}{3} = 5\sqrt{2}$. 

 Now from the sine area formula, the area of the quadrilateral is just$\frac{1}{2} \sin(\theta) ac + \frac{1}{2} \sin(\theta) bd + \frac{1}{2} \sin(180 - \theta) bc + \frac{1}{2} \sin(180 - \theta) ad = \frac{1}{2} \sin(\theta)(ac + bd + bc + ad) = \frac{1}{2} \sin(\theta)(a + b)(c + d) = \sin(\theta) \cdot a \cdot (c + d)$ where the last steps follow from $a = b$. Now $c + d = 3\sqrt{2} + 5\sqrt{2} = 8\sqrt{2}$. We need to find $a$ and $\sin(\theta)$.

 We can go back to equation (3) to find $a$. We have

 $a^2 + c^2 - 2ac\cos(\theta) = 100$

 Plugging in $c = 3\sqrt{2}$ and $\cos(\theta) = \frac{4c}{3a} = \frac{4 \cdot 3\sqrt{2}}{3a} = \frac{4\sqrt{2}}{a}$ we get

 $a^2 + 18 - 6\sqrt{2} \cdot a \cdot \frac{4\sqrt{2}}{a} = 100 \implies a^2 + 18 - 48 = 100 \implies a^2 = 130$

 Now $\sin(\theta) = \sqrt{1 - \frac{32}{a^2}}$ following the Pythagorean Identity. Now $a^2 = 130$ so we have

 $\sin(\theta) = \sqrt{1 - \frac{32}{130}} = \sqrt{\frac{98}{130}} = \frac{7\sqrt{2}}{\sqrt{130}}$. Hence the answer is 

 $\frac{7\sqrt{2}}{\sqrt{130}} \cdot 8\sqrt{2} \cdot \sqrt{130} = 7\sqrt{2} \cdot 8\sqrt{2} = 56 \cdot 2 = \boxed{112}$.

 



\textbf{Solution 13(cheese)}

The $AD = 2\sqrt{65}$ condition motivates us to find right triangles with integer sides and $2\sqrt{65}$ as the hypotenuse. Indeed, we have $2^2 + 16^2 = (2\sqrt{65})^2$. Furthermore, notice that $10$ is the hypotenuse of a $6-8-10$ right triangle. It follows that quadrilateral $ABCD$ is an $8 \times 22$ rectangle with triangles of sides $6-8-10$, $2-16-2\sqrt{65}$, and $6-8-10$ cut out from three of its corners. Thus, the area of $ABCD$ is $(8)(22) - \frac{1}{2}((6)(8) + (2)(16) + (6)(8)) = \boxed{112}$.

 This solution is very good for quick computational contests like AIME(I did NOT write this solution btw)~ilikemath247365

 


	\item (\textit{AIME II 2018}) Misha rolls a standard, fair six-sided die until she rolls 1-2-3 in that order on three consecutive rolls. The probability that she will roll the die an odd number of times is $\dfrac{m}{n}$ where $m$ and $n$ are relatively prime positive integers. Find $m+n$.


 
 \textbf{Answer: }647



\textbf{Solution 0}

Let $P_n$ be the probability of getting consecutive $1,2,3$ rolls in $n$ rolls and not rolling $1,2,3$ prior to the nth roll. 

 Let $x = P_3+P_5+...=1-(P_4+P_6+..)$. Following Solution 2, one can see that
 \[P_{n+1}=P_{n}-\frac{P_{n-2}}{6^3}\]for all positive integers $n \ge 5$. Summing for $n=5,7,...$ gives
 \[(1-x)-\frac{1}{6^3}=x-\frac{1}{6^3}-\frac{x}{6^3}\]
 \[\implies x = \frac{m}{n} = \frac{216}{431} \implies m+n=216+431= \boxed{647}\]

 



\textbf{Solution 1}

Let $P_{odd}=\frac{m}{n}$, with the subscript indicating an odd number of rolls. Then $P_{even}=1-\frac{m}{n}$.

 The ratio of $\frac{P_{odd}}{P_{even}}$ is just $\frac{P_{odd}}{1-P_{odd}}$.

 We see that $P_{odd}$ is the sum of $P_3$,$P_5$,$P_7$,... , while $P_{even}$ is the sum of $P_4$, $P_6$, $P_8$,... .

 $P_3$, the probability of getting rolls of 1-2-3 in exactly 3 rolls, is obviously $\frac{1}{216}$.

 We set this probability of $P_3$ aside, meaning we totally remove the chance of getting 1-2-3 in 3 rolls. Now the ratio of $P_4+P_6+P_8+...$ to $P_5+P_7+P_9+...$ should be equal to the ratio of $\frac{P_{odd}}{P_{even}}$, because in this case the 1st roll no longer matters, so we can disregard the very existence of it in counting how many times of rolls, and thus, 4 rolls, 6 rolls, 8 rolls... would become an odd number of rolls (while 5 rolls, 7 rolls, 9 rolls... would become even number of rolls).

 Notice $P_4+P_6+P_8+...=P_{even}$, and also $P_5+P_7+P_9+...=P_{odd}-P_3=P_{odd}-\frac{1}{216}$

 So we have $\frac{P_{even}}{P_{odd}-\frac{1}{216}}=\frac{P_{odd}}{P_{even}}$.

 Finally, we get $P_{odd}=\frac{m}{n}=\frac{216}{431}$. Therefore, $m+n = \boxed{647}$.

 



\textbf{Solution 2}

Call the probability you win on a certain toss $f_n$, where $n$ is the toss number.Obviously, since the sequence has length 3, $f_1=0$ and $f_2=0$.Additionally, $f_3=\left(\frac{1}{6}\right)^3$. We can call this value $x$, to keep our further equations looking clean.We can now write our general form for $f$ as $f_n=x\left(1-\sum_{i=1}^{n-3}f_i\right)$. This factors the probability of the last 3 rolls being 1-2-3, and the important probability that the sequence has not been rolled in the past (because then the game would already be over).Note that $\sum_{i=1}^{\infty}f_i=1$ since you'll win at some point.An intermediate step here is figuring out $f_n-f_{n+1}$. This is equal to $x\left(1-\sum_{i=1}^{n-3}f_i\right)-x\left(1-\sum_{i=1}^{n-2}f_i\right)=x\left(\sum_{i=1}^{n-2}f_i-\sum_{i=1}^{n-3}f_i\right)=xf_{n-2}$.Adding up all the differences, i.e. $\sum_{i=2}^{\infty}(f_{2n-1}-f_{2n})$ will give us the amount by which the odds probability exceeds the even probability. Since they sum to 1, that means the odds probability will be half of the difference above one-half. Subbing in our earlier result from the intermediate step, the odd probability is equal to $\frac{1}{2}+\frac{1}{2}x\sum_{i=2}^{\infty}f_{2n-3}$.Another way to find the odd probability is simply summing it up, which turns out to be $\sum_{i=1}^{\infty}f_{2n-1}$. Note the infinite sums in both expressions are equal; let's call it $P$. Equating them gives $\frac{1}{2}+\frac{1}{2}xP=P$, or $P=\frac{1}{2-x}$.Finally, substituting $x=\frac{1}{216}$, we find that $P=\frac{216}{431}$, giving us a final answer of $216 + 431 = \boxed{647}$.

 



\textbf{Solution 3}

Let $S(n)$ be the number of strings of length $n$ containing the digits $1$ through $6$ that do not contain the string $123$. Then we have $S(n) = 6 \cdot S(n-1) - S(n-3)$ because we can add any digit to end of a string with length $n-1$ but we have to subtract all the strings that end in $123$. We rewrite this as
 \begin{align*} S(n) &= 6 \cdot S(n-1) - S(n-3) \\ &= 6 \cdot (6 \cdot S(n-2) - S(n-4)) - (6 \cdot (S(n-4) - S(n-6)) \\ &= 36 \cdot S(n-2) - 12 \cdot S(n-4) + S(n-6) \end{align*}
 We wish to compute $P=\sum_{n=0}^\infty \frac{S(2n)}{6^{2n+3}}$ since the last three rolls are $123$ for the game to end. Summing over the recursion, we obtain
 \[\sum_{n=0}^\infty \frac{S(2n)}{6^{2n+3}} =36 \cdot \sum_{n=0}^\infty \frac{S(2n-2)}{6^{2n+3}} - 12 \cdot  \sum_{n=0}^\infty \frac{S(2n-4)}{6^{2n+3}}+ \sum_{n=0}^\infty \frac{S(2n-6)}{6^{2n+3}}\]Now shift the indices so that the inside term is the same:
 \begin{align*} \sum_{n=3}^\infty  \frac{S(2n)}{6^{2n+3}} &= \frac{36}{6^2} \cdot \sum_{n=2}^\infty \frac{S(2n)}{6^{2n+3}} - \frac{12}{6^4} \cdot \sum_{n=1}^\infty \frac{S(2n)}{6^{2n+3}} + \frac{1}{6^6} \cdot \sum_{n=0}^\infty \frac{S(2n)}{6^{2n+3}}  \\  \left(P - \frac{S(0)}{6^3} -  \frac{S(2)}{6^5} -\frac{S(4)}{6^7}  \right) &= \frac{36}{6^2} \cdot \left( P -   \frac{S(0)}{6^3} -  \frac{S(2)}{6^5}\right) - \frac{12}{6^4} \cdot \left(  P - \frac{S(0)}{6^3} \right) + \frac{1}{6^6} \cdot P  \end{align*}
Note that $S(0) = 1, S(2) = 36$ and $S(4) = 6^4 - 2 \cdot 6 = 1284$. Therefore,
 \begin{align*} \left(P - \frac{1}{6^3} -  \frac{36}{6^5} -\frac{1284}{6^7}  \right) = \frac{36}{6^2} \cdot \left( P -   \frac{1}{6^3} -  \frac{36}{6^5}\right) - \frac{12}{6^4} \cdot \left(  P - \frac{1}{6^3} \right) + \frac{1}{6^6} \cdot P   \end{align*}
Solving for $P$, we obtain $P = \frac{216}{431} \implies m+n = \boxed{647}$. 

 -Vfire

 



\textbf{Solution 4}

Let $A=\frac{1}{6} \begin{bmatrix} 5 & 1 & 0 & 0 \\ 4 & 1 & 1 & 0 \\ 4 & 1 & 0 & 1 \\ 0 & 0 & 0 & 0 \\ \end{bmatrix}$. $A$ is a transition matrix for the prefix of 1-2-3 matched so far. The state corresponding to a complete match has no outgoing probability mass. The probability that we roll the dice exactly $k$ times is $(A^k)_{1,4}$. Thus the probability that we roll the dice an odd number of times is $1-\left(\sum_{k=0}^\infty A^{2k}\right)_{1,4} = 1-\left((I - A^2)^{-1}\right)_{1,4} = \frac{216}{431}$. Thus the answer is $216+431=\boxed{647}$.

 



\textbf{Solution 5 (quick cheat)}

Consider it as a contest of Odd and Even. Let $P_o$ and $P_e$ be probability that Odd and Even wins, respectively. If we consider every 3 rolls as an atomic action, then we can have a simple solution. If the rolls is 1-2-3, Odd wins; otherwise, Odd and Even switch the odds of winning. Therefore, we have
 \[P_o = \frac{1}{216} + \frac{215}{216}P_e\]Plug in $P_e = 1 - P_o$ and we can easily solve for $Po=\frac{216}{431}$. 

 $\boxed{647}$.

 Of course this is not a rigorous solution. I think it works because the requirement is a strict sequence of pure random events.

 -Mathdummy

 



\textbf{Solution 6 (Markov Chain)}


            \begin{enumerate}[label=(\roman*)]
                \item <i>Main article: \href{https://artofproblemsolving.com/wiki/index.php?title=Markov\_Chains}{Markov Chains}</i>
            \end{enumerate}
            Let $P_o$, $P_e$ be the winning probabilities respectively. We call Odd "in position" when a new sequence of 1-2-3 starts at odd position, and likewise, call Even is "in position" when a new sequence starts at even position. 

 Now, consider the situation when the first roll is $1.$ The conditional probability for Odd or Even to eventually win out depends on whose is in position. So let's denote by $P_o(1), P_e(1)$ the probabilities of Odd and Even winning out, respectively, both when Odd is in position. Remember that the probabilities simply switch if Even is in position. Similarly, after 1-2 is rolled, we denote by $P_o(2), P_e(2)$ the conditional probabilities of Odd and Even winning out, when Odd is in position.

 Consider the first roll. If it's not a 1, the sequence restarts, but Even is now in position; if it's a 1, then Odd's winning probability becomes $P_o(1)$. So,
 \[P_o = \frac{1}{6}P_o(1) + \frac{5}{6}P_e\]In the next roll, there are 3 outcomes. If the roll is 2, then Odd's winning probability becomes $P_o(2)$; if the roll is 1, then we stay in the sequence, but Even is now in position, so the probability of Odd winning now becomes $P_e(1)$; if the rolls is any other number, then the sequence restarts, and Odd is still in position. So,
 \[P_o(1) = \frac{1}{6}P_o(2) + \frac{1}{6}P_e(1) + \frac{4}{6}P_o\]In the next roll after a 1-2 sequence, there are 3 outcomes. If the roll is a 3, Odd wins; if it's a 1, we go back to the state when 1 is just rolled, and Odd is in position; if it's any other number, then the sequence restarts, and Even is in position. So,
 \[P_o(2) = \frac{1}{6} + \frac{1}{6}P_o(1) + \frac{4}{6}P_e\]

 Plug in $P_e = 1-P_o$ and $P_e(1) = 1 - P_o(1)$, we have a 3-equation linear system which is not hard to solve. The final answer is $Po=\frac{216}{431}$. $\boxed{647}$. (We want P\_o because if it starts odd, it will also end odd)

 -Mathdummy

 \textbf{Markov Chain Diagram} 
                    \begin{figure}[H]
                        \centering
                        \includegraphics[width=0.5\textwidth]{Solution/AIME_II_2018_13_Solution_Figure_1.png}
                    \end{figure}
                    

 \textbf{mathboy282}

 $\LaTeX$ done by \href{https://artofproblemsolving.com/wiki/index.php?title=User:Crazyvideogamez}{CrazyVideoGamez}

 



\textbf{Solution 7(Blocks)}

Take a block of any three possible rolls. You have $6^3$ different possibilities for a block and $6^3 - 1$ different possibilities to have a block without 1-2-3. This gives two probabilities of $\frac{6^3-1}{6^3}$ and $\frac{1}{6^3}$, respectively. To get an odd number of total rolls, we need an even number of blocks without a 1-2-3. Then, we can sum up the probabilities up to infinity to see the total probability of getting an even number of rolls.
 \[\frac{1}{6^3}\sum_{k=0}^{\infty}\left(\frac{6^3-1}{6^3}\right)^{2k}\]
 \[=\frac{\frac{1}{6^3}}{1-\left(\frac{6^3-1}{6^3}\right)^2}\]Using difference of squares on the bottom half we can reduce it to $\left(1-\frac{6^3-1}{6^3}\right)\left(1+\frac{6^3-1}{6^3}\right)=\frac{2\times 6^3 -1}{(6^3)^2}$. Plugging this back in to the equation we get- 
 \[\frac{\frac{1}{6^3}}{\frac{2\times 6^3 -1}{(6^3)^2}} = \frac{6^3}{2\times 6^3 - 1} = \frac{216}{431}\]Then, our answer is $\boxed{647}$-Mathiscool109

 



\textbf{Solution 8 (Recursion)}

Let the $P_o$ be the probability that Misha will roll the die an odd number of times, $P_e$ be the probability that Misha will roll the die an even number of times.

 Let $P_n$ be the probability Misha ends after rolling $n$ number of times.

 Ending at $n$ has two requirements: The last three rolls are $1,2,3$, and there can't exist any $1,2,3$ in previous rolls, which causes Misha end at a $k < n$.

 $P_n = (\frac{1}{6})^3 \cdot (1-\sum\limits_{k=1}\limits^{n-3} P_k)$

 Substitute the values for

 $\begin{cases} P_{2n+1} = \frac{1}{216} (1-\sum\limits_{k=1}\limits^{2n-2} P_k) \\ P_{2n} = \frac{1}{216} (1-\sum\limits_{k=1}\limits^{2n-3} P_k) \end{cases}$

 
$P_{2n+1}-P_{2n} = \frac{1}{216}(-P_{2n-2})$

 $\sum\limits_{n=2}\limits^{\infty} P_{2n+1} - \sum\limits_{n=2}\limits^{\infty} P_{2n} = -\frac{1}{216}(\sum\limits_{n=2}\limits^{\infty} P_{2n-2})$

 $(P_o - P_3) - P_e = -\frac{1}{216} \cdot P_e$

 It is easy to see that $P_3 = -(\frac{1}{6})^3 = \frac{1}{216}$ with only one case of $1,2,3$. Also, $P_{o}$ and $P_{e}$ should be complementary values. Therefore we get two equations for two unknown variables:

 
 \[\begin{cases} P_o - \frac{1}{216} = \frac{215}{216} P_e \\ P_o + P_e = 1 \end{cases}\]

 
So $\frac{431}{216} \cdot P_e = \frac{215}{216}$, and $P_e = \frac{215}{431}$

 $P_o = 1-\frac{215}{431} = \frac{216}{431}$

 $216+431=\boxed{647}$

 -cassphe

 


	\item (\textit{AIME II 2018}) The incircle $\omega$ of triangle $ABC$ is tangent to $\overline{BC}$ at $X$. Let $Y \neq X$ be the other intersection of $\overline{AX}$ with $\omega$. Points $P$ and $Q$ lie on $\overline{AB}$ and $\overline{AC}$, respectively, so that $\overline{PQ}$ is tangent to $\omega$ at $Y$. Assume that $AP = 3$, $PB = 4$, $AC = 8$, and $AQ = \dfrac{m}{n}$, where $m$ and $n$ are relatively prime positive integers. Find $m+n$.


 
 \textbf{Answer: }227



\textbf{Solution 1}

Let the sides $\overline{AB}$ and $\overline{AC}$ be tangent to $\omega$ at $Z$ and $W$, respectively. Let $\alpha = \angle BAX$ and $\beta = \angle AXC$. Because $\overline{PQ}$ and $\overline{BC}$ are both tangent to $\omega$ and $\angle YXC$ and $\angle QYX$ subtend the same arc of $\omega$, it follows that $\angle AYP = \angle QYX = \angle YXC = \beta$. By equal tangents, $PZ = PY$. Applying the Law of Sines to $\triangle APY$ yields 
 \[\frac{AZ}{AP} = 1 + \frac{ZP}{AP} = 1 + \frac{PY}{AP} = 1 + \frac{\sin\alpha}{\sin\beta}.\]Similarly, applying the Law of Sines to $\triangle ABX$ gives 
 \[\frac{AZ}{AB} = 1 - \frac{BZ}{AB} = 1 - \frac{BX}{AB} = 1 - \frac{\sin\alpha}{\sin\beta}.\]It follows that 
 \[2 = \frac{AZ}{AP} + \frac{AZ}{AB} = \frac{AZ}3 + \frac{AZ}7,\]implying $AZ = \frac{21}5$. Applying the same argument to $\triangle AQY$ yields 
 \[2 = \frac{AW}{AQ} + \frac{AW}{AC} = \frac{AZ}{AQ} + \frac{AZ}{AC} = \frac{21}5\left(\frac{1}{AQ} + \frac 18\right),\]from which $AQ = \frac{168}{59}$. The requested sum is $168 + 59 = \boxed{227}$.

 



\textbf{Solution 2 (Projective)}

Let the incircle of $ABC$ be tangent to $AB$ and $AC$ at $Z$ and $W$. By Brianchon's theorem on tangential hexagons $QWCBZP$ and $PYQCXB$, we know that $ZW,CP,BQ$ and $XY$ are concurrent at a point $O$. Let $PQ \cap BC = M$. Then by La Hire's $A$ lies on the polar of $M$ so $M$ lies on the polar of $A$. Therefore, $ZW$ also passes through $M$. Then projecting through $M$, we have
 \[-1 = (A,O;Y,X) \stackrel{M}{=} (A,Z;P,B) \stackrel{M}{=} (A,W;Q,C).\]Therefore, $\frac{AP \cdot ZB}{MP \cdot AB} = 1 \implies \frac{3 \cdot ZB}{ZP \cdot 7} = 1$. Since $ZB+ZP=4$ we know that $ZP = \frac{6}{5}$ and $ZB = \frac{14}{5}$. Therefore, $AW = AZ = \frac{21}{5}$ and $WC = 8 - \frac{21}{5} = \frac{19}{5}$. Since $(A,W;Q,C) = -1$, we also have $\frac{AQ \cdot WC}{NQ \cdot AC} = 1 \implies \frac{AQ \cdot \frac{19}{5}}{(\frac{21}{5} - AQ) \cdot 8} = 1$. Solving for $AQ$, we obtain $AQ = \frac{168}{59} \implies m+n = \boxed{227}$.😃-Vfire

 



\textbf{Solution 3  (Combination of Law of Sine and Law of Cosine)}

Let the center of the incircle of $\triangle ABC$ be $O$. Link $OY$ and $OX$. Then we have $\angle OYP=\angle OXB=90^{\circ}$

 $\because$ $OY=OX$

 $\therefore$ $\angle OYX=\angle OXY$

 $\therefore$ $\angle PYX=\angle YXB$

 $\therefore$ $\sin \angle PYX=\sin \angle YXB=\sin \angle YXC=\sin \angle PYA$

 Let the incircle of $ABC$ be tangent to $AB$ and $AC$ at $M$ and $N$, let $MP=YP=x$ and $NQ=YQ=y$.

 Use Law of Sine in $\triangle APY$ and $\triangle AXB$, we have

 $\frac{\sin \angle PAY}{PY}=\frac{\sin \angle PYA}{PA}$

 $\frac{\sin \angle BAX}{BX}=\frac{\sin \angle AXB}{AB}$

 therefore we have

 $\frac{3}{x}=\frac{7}{4-x}$

 Solve this equation, we have $x=\frac{6}{5}$

 As a result, $MB=4-x=\frac{14}{5}=BX$, $AM=x+3=\frac{21}{5}=AN$, $NC=8-AN=\frac{19}{5}=XC$, $AQ=\frac{21}{5}-y$, $PQ=\frac{6}{5}+y$

 So, $BC=\frac{14}{5}+\frac{19}{5}=\frac{33}{5}$

 Use Law of Cosine in $\triangle BAC$ and $\triangle PAQ$, we have

 $\cos \angle BAC=\frac{AB^2+AC^2-BC^2}{2\cdot AB\cdot AC}=\frac{7^2+8^2-{(\frac{33}{5})}^2}{2\cdot 7\cdot 8}$

 $\cos \angle PAQ=\frac{AP^2+AQ^2-PQ^2}{2\cdot AP\cdot AQ}=\frac{3^2+{(\frac{21}{5}-y)}^2-{(\frac{6}{5}+y)}^2}{2\cdot {(\frac{21}{5}-y)}\cdot 3}$

 And we have

 $\cos \angle BAC=\cos \angle PAQ$

 So

 $\frac{7^2+8^2-{(\frac{33}{5})}^2}{2\cdot 7\cdot 8}=\frac{3^2+{(\frac{21}{5}-y)}^2-{(\frac{6}{5}+y)}^2}{2\cdot {(\frac{21}{5}-y)}\cdot 3}$

 Solve this equation, we have $y=\frac{399}{295}=QN$

 As a result, $AQ=AN-QN=\frac{21}{5}-\frac{399}{295}=\frac{168}{59}$

 So, the final answer of this question is $168+59=\boxed {227}$

 



\textbf{Solution 4 (Projective geometry)}


                    \begin{figure}[H]
                        \centering
                        \includegraphics[width=0.5\textwidth]{Solution/AIME_II_2018_14_Solution_Figure_1.png}
                    \end{figure}
                    <i>\textbf{Claim}</i>

 Let the sides $\overline{AB}$ and $\overline{AC}$ be tangent to $\omega$ at $M$ and $N$, respectively. Then lines $PQ, MN,$ and $BC$ are concurrent and lines $PC, MN, AX,$ and $BQ$ are concurrent.

 <i>\textbf{Proof}</i>

 Let $E$ be point of crossing $AX$ and $BQ.$ We make projective transformation such that circle $\omega$ maps into the circle and point $E$ maps into the center of new circle point $I.$ We denote images using notification $X \rightarrow X'.$

 
                    \begin{figure}[H]
                        \centering
                        \includegraphics[width=0.5\textwidth]{Solution/AIME_II_2018_14_Solution_Figure_2.png}
                    \end{figure}
                    $BCQP$ maps into $B'C'Q'P'$, so lines $B'Q'$ and  $A'X'$ be the diameters.This implies $P'Q'||B'C', \angle B'P'Q' = \angle B'C'Q' = 90^\circ \implies B'C'Q'P'$ be a square.

 Therefore $M'N'$ be the diameter $\implies P'C', B'Q',$ be diagonals of the square. $M'N'$ and  $X'Y'$ be midlines  which crossing in the center $I.$ Therefore lines $PC, MN, AX,$ and $BQ$ are concurrent.

 Lines $P'Q'||M'N' ||B'C' \implies PQ, MN$ and $BC$ are concurrent.

 <i>\textbf{Solution}</i>The cross-ratio associated with a list of four collinear points $A,P,M,D$ is defined as 
 \[(A,P;M,B)={\frac {AP\cdot MB}{AB\cdot PM}}.\]The cross-ratio be <i>\textbf{projective invariant}</i> of a quadruple of collinear points, so

 
                    \begin{figure}[H]
                        \centering
                        \includegraphics[width=0.5\textwidth]{Solution/AIME_II_2018_14_Solution_Figure_3.png}
                    \end{figure}
                    
 \[(A,P; M,B)  = {\frac {A'P'\cdot M'B'}{A'B'\cdot P'M'}} = \frac {M'B'}{P'M'} = 1.\] 
 \[(A,P; M,B)={\frac {3\cdot (7 - AM)}{7\cdot (AM - 3)}} = 1 \implies AM = \frac {21}{5} \implies AN = AM =  \frac {21}{5}.\]

 
                    \begin{figure}[H]
                        \centering
                        \includegraphics[width=0.5\textwidth]{Solution/AIME_II_2018_14_Solution_Figure_4.png}
                    \end{figure}
                    
 \[(A,Q;N,C)={\frac {AQ\cdot NC}{AC\cdot QN}} = \frac {AQ\cdot (AC- AN)}{AC\cdot (AN-AQ)} = 1.\] 
 \[AQ \cdot (8 - \frac{21}{5}) = 8 \cdot (\frac{21}{5} – AQ) \implies  AQ = \frac{168}{59}.\]

 For visuals only, I will show how one can find the perceptor $D$ and the image’s plane.$E_0$ is image of inversion $E$ with respect $\omega.$$UW$ is the diameter of $\omega, E,E_0,U,W$ are collinear.$DU \perp \omega, DE_0\perp WD, UV \perp WD, UV$ is diameter of $\omega'$.

 Plane of images is perpendicular to $WD.$

 Last diagram shows the result of transformation. Transformation is possible. The end.

 \textbf{vladimir.shelomovskii@gmail.com, vvsss}

 



\textbf{Solution 5}

Firstly, assume $PY=x=PZ, ZB=4-x=BX, AZ=AW=3+x, CW=CX=5-x, QY=QW=y, AQ=3+x-y$

 By tangency, we have $\angle{PYZ}=\angle{PZY}=\angle{YXZ}; \angle{BZX}=\angle{BXZ}=\angle{ZYX}; \angle{PYX}=\angle{BXY}$

 Similar reason yields $\angle{QYX}=\angle{CXY}$. Apply Law of sines

 We have $\frac{3}{\sin{\angle{PYA}}}=\frac{x}{\sin{\angle{PAY}}}, \frac{7}{\sin{\angle{AXB}}}=\frac{4-x}{\sin{\angle{BAX}}}$ Since $\angle{PYA}=180-\angle{BXA}$ so their sine values would be the same. Solve this system and we have $\frac{3}{7}=\frac{x}{4-x}, x=\frac{6}{5}$

 Apply the same process in $\triangle{AQY}, \triangle{AXC}$, we have $\frac{3+x-y}{8}=\frac{y}{5-x}, y=\frac{399}{295}$

 The desired length is $3+x-y=\frac{168}{59}\implies \boxed{227}$

 



\textbf{Solution 6(Algebra Bash)}

Let $\angle PAC = \theta$. Let $PY = a$ and $YQ = c$. By PoP and Stewart's Theorem

 $(4 - a)(9 - 2a)(5 - a) + AX^2 (9 - 2a) = 64(4 - a) + 49(5 - a)$ where we will let $AX = u$ for simplicity. Now let $AY = x$. By LoC we have

 $9 + x^2 - 6x\cos(\theta) = a^2$

 By LoC we also have

 $49 + AX^2 - 14 AX \cos(\theta) = (4 - a)^2$

 Finally, by PoP, we have

 $x \cdot AX = (a + 3)^2$

 Solving this system of equations, we get

 $a = \frac{6}{5}, x = \frac{63 \cdot \sqrt{2761}}{1255}$. Note that by PoP, $AQ = a + 3 - c$ by using the Power of Point of tangents that are equal. Finally, by Stewart's Theorem on $\triangle APQ$ we have

 $a(a + c)c + x^2 (a + c) = (a + 3 - c)^2 a + 9c$ in which after plugging in the values we found for $a$ and $x$ we get $c = \frac{399}{295}$. Finally, $a + 3 - c = \frac{168}{59} \implies \boxed{227}$.

 \textbf{Proof of why (A, E; Y, X) is a harmonic bundle without projective geo} Consider a circle and let $A$ be a point outside the circle. Let the tangents from $A$ intersect the circle at points $M$ and $N$. That's exactly what's happening in this problem. Now, the line $AX$ is a line that cuts the circle at point $Y$. Let $AX \cap MN = E$. 

 Lemma: Points $A, Y, E, X$ form a harmonic bundle or $(A, E; Y, X) = -1 \implies \frac{YA}{YE} = \frac{XA}{XE}$ ignoring directed lengths.

 Proof: $\triangle AMN$ is isosceles since $AM, AN$ are tangents. So we can apply Stewart's Theorem to get

 $EN \cdot MN \cdot EM + AE^2 \cdot MN = EN \cdot AM^2 + EM \cdot AN^2$ but $AM = AN$ so we have

 $EN \cdot MN \cdot EM + AE^2 \cdot MN = AM^2(EN + EM) = AM^2(MN) \implies EN \cdot EM + AE^2 = AM^2$

 Now, note that $Pow(E) = EM \cdot EN = EY \cdot EX \implies AM^2 = EY \cdot EX + AE^2 = AE^2 + (AE - AY)(AX - AE) = AE^2 + AE \cdot AX - AE^2 - AY \cdot AX + AY \cdot AE = AE \cdot AX - AY \cdot AX + AY \cdot AE$ however note that $Pow(A) = AM^2 = AY \cdot AX$ so $AY \cdot AX = AE \cdot AX - AY \cdot AX + AY \cdot AE \implies 2 \cdot AX \cdot AY = AE(AX + AY) \implies \frac{2}{AE} = \frac{1}{AX} + \frac{1}{AY}$ which means that $AE$ is the harmonic mean of $AX$ and $AY$ and therefore $(A, E; Y, X)$ is a harmonic bundle as desired. $\square$

 


	\item (\textit{AIME II 2018}) Find the number of functions $f$ from $\{0, 1, 2, 3, 4, 5, 6\}$ to the integers such that $f(0) = 0$, $f(6) = 12$, and


 
 \[|x - y|  \leq  |f(x) - f(y)|  \leq  3|x - y|\]


 for all $x$ and $y$ in $\{0, 1, 2, 3, 4, 5, 6\}$.


 
 \textbf{Answer: }185



\textbf{Solution 1}

First, suppose $x = y + 1$. Then, the inequality becomes $1 \leq |f(y + 1) - f(y)| \leq 3$. In other words, the (positive) difference between consecutive function values is $1$, $2$, or $3$. Let $d_k := f(k) - f(k - 1)$. Note that

 $f(x) - f(y) = \sum_{k=y+1}^{x} d_k.$

 Thus, $\sum_{k=1}^{6} d_k = f(6) - f(0) = 12$. Note that at most one value of $d_k$ in $\left\{d_1, d_2, d_3, d_4, d_5, d_6\right\}$ can be negative. This is because the maximum value of $d_1 + d_2 + d_3 + d_4 + d_5 + d_6$ would be $3 + 3 + 3 + 3 - 1 - 1 = 10$ if more than one value of $d_k$ is negative. Plugging $x = y + 2$ into the original inequality yields $2 \leq |f(y + 2) - f(y)| \leq 6$, which becomes $2 \leq \left|d_{y+2} + d_{y+1}\right| \leq 6$. The only way for $d_{y+2} + d_{y+1}$ to be negative while satisfying this inequality is for $\left(d_{y+2}, d_{y+1}\right)$ to equal $(1, -3)$ or $(-3, 1)$. However, this forces $d_1 + d_2 + d_3 + d_4 + d_5 + d_6 \leq 3 + 3 + 3 + 3 + 1 - 3 = 10 < 12$, which is disallowed. Hence, we conclude that the following stronger inequality,

 $2 \leq d_{y + 2} + d_{y + 1} \leq 6,$

 is always true. We now have two cases of functions to count. For future reference let $D$ be the (ordered) sequence $\left\{d_1, d_2, d_3, d_4, d_5, d_6\right\}$.

 Case 1: There is exactly one instance of $d_k = -1$.

 By the "stronger" inequality above, $d_{k-1} = 3$ if $k > 1$, and $d_{k+1} = 3$ if $k < 6$. If $2 \leq k \leq 5$, then $D$ contains the subsequence $\left\{3, -1, 3\right\}$, and the other three $d$-values sum to $7$, so they are either $3$, $3$, and $1$ in some order, or they are $2$, $2$, and $3$ in some order. Thus, each $k \in \{2, 3, 4, 5\}$ for which $d_k = -1$ produces $6$ sequences $D$. If $k = 1$ or $k = 6$, then $D$ begins with $\{-1, 3\}$ or ends with $\{3, -1\}$, respectively. Either way, the remaining four $d$-values sum to $10$, so they can be any permutation of $\{3, 3, 2, 2\}$ (six permutations) or $\{3, 3, 3, 1\}$ (four permutations). Each of these vaues of $k$ yields $6 + 4 = 10$ sequences, so our total count for Case 1 is $4 \cdot 6 + 2 \cdot 10 = 44$.

 Case 2: All values of $d$ are positive.

 Then, $D$ is a permutation of $\{3, 3, 3, 1, 1, 1\}$, $\{3, 3, 2, 2, 1, 1\}$, $\{3, 2, 2, 2, 2, 1\}$, or $\{2, 2, 2, 2, 2, 2\}$.The number of ways to permute three $3$s and three $1$s is $\frac{6!}{3! \cdot 3!} = 20$.The number of ways to permute two $3$s, two $2$s, and two $1$s is $\frac{6!}{2! \cdot 2! \cdot 2!} = 90$.The number of ways to permute one $3$, four $2$s, and one $1$ is $\frac{6!}{1! \cdot 4! \cdot 1!} = 30$.Finally, there is obviously only $1$ way to permute six $2$s. Our total count for Case 2 is $20 + 90 + 30 + 1 = 141$.

 To complete the justification that all of these $44 + 141 = \boxed{185}$ cases satisfy the original inequality, we leverage the fact that $\{f(x)\}_{x=0}^{6}$ is either monotonically increasing (Case 2) or the union of two monotonically increasing subsequences (Case 1). Consider any monotonically increasing subsequence that starts at $x = a$ and ends at $x = b$. For $a < k \leq b$, $d_k$ will be positive, allowing us to remove the absolute value bars from the original inequality:

 $x - y \leq f(x) - f(y) \leq 3(x - y).$

 Now, the inequality is transitive; supposing $a \leq x_3 < x_2 < x_1 \leq b$, if the inequality is satisfied at $(x, y) = \left(x_1, x_2\right)$ and at $(x, y) = \left(x_2, x_3\right)$, then it is also satisfied at $(x, y) = \left(x_1, x_3\right)$. If we ever have a decreasing part where $f(x + 1) < f(x)$, then we can use some variant of the inequality $2 \leq f(x + 1) - f(x - 1) \leq 6$, which we derived earlier. This is a specific case of $x - y \leq f(x) - f(y) \leq 3(x - y)$, so we can finish off the argument by invoking transitivity.    

 -kgator

 



\textbf{Solution 2}

Based on the conditions for $d_k$ obtained in Solution 1, followed by some well-known stars & bars techniques to reduce the case work. To recap, the conditions for $d_k = f(k) - f(k-1)$ are:
 \[\sum_{k=1}^6 d_k = 12\]
 \[d_k = -1, \hspace{0.1cm} \mbox{or} \hspace{0.2cm}  1 \leq d_k \leq 3\]
 \[d_k + d_{k+1} \geq 2\]

 Case 1: $d_k = -1$ for some $k$. Since $d_i + d_{i+1} \geq 2$, $d_{k+1}$ and $d_{k-1}$ must both be 3. 

 Case 1a: $d_1 = -1$ or $d_6 = -1$. In this case, we have $d_2=3$ or $d_5=3$, respectively. So the remaining four $d_i$ values add up to 10 with each value ranging between 1 and 3. This stars and bars problem is equivalent to $a+b+c+d =2$ with $a,b,c,d$ between 0 and 2, so the number of possibilities is ${5 \choose 3}$ each.

 Case 1b: $d_k = -1$ for $k=2,3,4,5$. In this case, $d_{k+1}=d_{k-1} =3$. So the remaining 3 $d_i$ values add up to 7 with each value ranging between 1 and 3. This stars and bars problem is equivalent to $a+b+c =2$ with $a,b,c,d$ between 0 and 2, so the number of possibilities is ${4 \choose 2}$ each.

 So the total for Case 1 is $2{5 \choose 3} + 4 {4 \choose 2} = 44$.

 Case 2: $1 \leq d_k \leq 3$ for all $k$. $\sum_{k=1}^{6} d_k = 12$. Let $g_k = 3-d_k$ we get $\sum_{k=1}^{6} g_k = 6$ for $0 \leq g_k \leq 2$. This is a stars and bars problem with maximum capacity and the solution is given by applying PIE:
 \[{11 \choose 5} - 6{8 \choose 5} + {6 \choose 2} = 141\].Now the final answer is $44 + 141 = 185$.

 -mathdummy

 


\end{enumerate}